\documentclass[11pt]{article}
\usepackage[T1]{fontenc}
\usepackage[utf8]{inputenc}
\usepackage{lmodern}
\usepackage[a4paper,margin=26mm,headheight=14pt]{geometry}
\usepackage{amsmath,amssymb,amsthm,mathtools}
\usepackage{booktabs,longtable,array}
\usepackage{microtype}
\usepackage{xcolor}
\usepackage{enumitem}
\usepackage{fancyhdr}
\usepackage{hyperref}
\hypersetup{colorlinks=true,linkcolor=blue!45!black,citecolor=blue!45!black,urlcolor=blue!45!black,
 pdftitle={Unconditional all-orders asymptotics for two OEIS permutation sequences},
 pdfauthor={Research report prepared for Vladimir Reshetnikov}}
\usepackage[nameinlink,noabbrev]{cleveref}
\numberwithin{equation}{section}
\newtheorem{theorem}{Theorem}[section]
\newtheorem{proposition}[theorem]{Proposition}
\newtheorem{lemma}[theorem]{Lemma}
\newtheorem{corollary}[theorem]{Corollary}
\theoremstyle{definition}
\newtheorem{definition}[theorem]{Definition}
\newtheorem{conjecture}[theorem]{Conjecture}
\theoremstyle{remark}
\newtheorem{remark}[theorem]{Remark}
\newcommand{\E}{\mathbb E}
\newcommand{\Pp}{\mathbb P}
\newcommand{\Z}{\mathbb Z}
\newcommand{\Q}{\mathbb Q}
\newcommand{\N}{\mathbb N}
\newcommand{\fall}[2]{{#1}^{\underline{#2}}}
\newcommand{\rise}[2]{{#1}^{\overline{#2}}}
\newcommand{\bfac}{\mathbf b!}
\newcommand{\betafac}{\boldsymbol\beta!}
\newcommand{\coef}[2]{[#1^{#2}]}
\newcommand{\TV}{d_{\mathrm{TV}}}
\newcommand{\ind}{\mathbf 1}
\newcommand{\Aone}{\textnormal{A189281}}
\newcommand{\Atwo}{\textnormal{A110128}}
\newcommand{\bigO}{\mathrm O}
\newcommand{\smallo}{\mathrm o}
\setlist{itemsep=3pt,topsep=5pt}
\setlength{\emergencystretch}{2em}
\pagestyle{fancy}
\fancyhf{}
\fancyhead[L]{\small Fixed-offset permutation avoidance}
\fancyhead[R]{\small Research report}
\fancyfoot[C]{\thepage}
\title{\textbf{Unconditional all-orders asymptotics\\for two OEIS permutation sequences}\\[5pt]
\large Stable path forests, Poisson expansions, and index inversion}
\author{Prepared for Vladimir Reshetnikov}
\date{1 October 2026}

\begin{document}
\maketitle
\begin{abstract}
The OEIS entries A189281 and A110128 count permutations avoiding a prescribed
signed or absolute difference between entries two positions apart. Their displayed
higher asymptotic expansions are associated with conjectured recurrences. We give
a direct, recurrence-free proof of those expansions, together with a finite exact
formula for every coefficient. More generally, for any fixed positive offsets
$r,s$, the complete probability generating function of the number of violations
has an all-orders expansion consisting of $e^{\theta u}$ times polynomial
corrections, where $\theta=1$ in the oriented case and $\theta=2$ in the absolute
case. The remainder is uniform on compact complex $u$-sets.

The structural ingredient is an exact stabilization theorem for tilings of long
path forests: low factorial moments depend on the numbers of paths, not their
individual lengths. We also prove a quantitative total-variation universality
bound, a factorial denominator bound for all asymptotic coefficients, and an
all-orders Lambert-$W$ index inversion. Exact rational computations recover the
OEIS coefficients through order ten and extend both expansions through order
sixteen. Independent permutation enumeration and Bonferroni certificates accompany
the proofs. The specific conjectured recurrences and the stronger all-orders
integrality suggested by A189281 are not assumed or claimed proved. A concrete
rational-function conjecture is proposed as a route to the latter question.
\end{abstract}

\paragraph{Status.}
This is a proof-focused research draft, not a claim of peer review or a
Lean-checked formalization. The theorems below have mathematical proofs; the
computations are reproducible exact checks. Historical priority beyond the
sources reviewed here remains to be assessed. The results establish the
algebraic asymptotic sector, not a complete exponentially improved transseries.

\paragraph{Editorial note (1 October 2026, batch 73O2).}
This report is manuscript~60 of batch~73 of ProveIt's incoming-reports intake,
placed in the research-report collection in commit \texttt{6e193dd4f}. It is
AI-assisted, unrefereed and not formalized. The text is the delivered
manuscript; the write added only the label prefix \texttt{spf:}, the bracketed
notes marked ``[Added 1 October 2026, batch 73O2]'',
Remarks~\ref{spf:rem:offsets-one-one} and~\ref{spf:rem:inverse-apparatus}, three bibliography
entries, and \cref{spf:app:provenance}, which records the provenance and every
editorial change. The rational-collapse statement of \cref{spf:sec:collapse}
remains a conjecture.

\paragraph{Editorial note (1 October 2026, batch 74).}
Since batch~74 the report has two Parts. Part~I is manuscript~60, as
described above. Part~II (\cref{fga:sec:front}) adds manuscript~03 of
batch~74, an independent second derivation of Part~I's main theorems, and
prints only what it adds: uniformity over forest shapes, comparison bounds,
closed forms of $c_3(r,s)$, and proved integrality for the oriented one-path
family. Its labels carry the prefix \texttt{fga:}.

\newpage
\tableofcontents
\newpage

\part{Unconditional all-orders asymptotics for two OEIS permutation sequences}

\section{The target and the precise contribution}
\label{spf:sec:target}

For a uniformly random $\pi\in S_n$, define, for fixed $r,s\ge1$,
\begin{align}
 X^{(1)}_{r,s,n}
   &=\sum_{i=1}^{n-r}\ind\{\pi(i+r)-\pi(i)=s\},\\
 X^{(2)}_{r,s,n}
   &=\sum_{i=1}^{n-r}\ind\{|\pi(i+r)-\pi(i)|=s\}.
\end{align}
An empty sum is zero. Write
\begin{equation}
 A^{(\theta)}_{r,s}(n)=n!\,\Pp(X^{(\theta)}_{r,s,n}=0),
 \qquad \theta\in\{1,2\}.
 \label{spf:eq:counts}
\end{equation}
In this notation,
\begin{equation}
 \Aone(n)=A^{(1)}_{2,2}(n),\qquad
 \Atwo(n)=A^{(2)}_{2,2}(n).
\end{equation}
These identifications and the sequence definitions are recorded in the OEIS
entries~\cite{oeis189281,oeis110128}.

\subsection{What the sources do, and do not, already establish}

The A189281 entry lists an order-eight, degree-eleven recurrence as conjectural,
and describes its displayed terms through $n^{-10}$ as consequences of that
recurrence. It also records independent dynamic-programming verification of
sequence values through $n=300$~\cite{oeis189281}. The recurrence-guessing work of
Kauers and Koutschan is an important source for this problem~\cite{kauers}.
For A110128, the entry lists a guessed order-24, degree-64 operator and attributes
its higher asymptotic terms to that operator~\cite{oeis110128}.

The common leading limits are not new. Tauraso proved, in the diagonal absolute
case, the first correction
\begin{equation}
 \frac{A^{(2)}_{r,r}(n)}{n!}
 =e^{-2}\left(1+\frac{4(r-1)}n+\bigO(n^{-2})\right)
 \label{spf:eq:tauraso}
\end{equation}
and gave exact counting formulas~\cite{tauraso}. Spahn and Zeilberger developed
an exact matching-of-tilings formula for both offset families, explained its
computational use, and discussed Matsuo's alternative approach~\cite{spahn}.
We use their combinatorial framework, with a self-contained proof, rather than
claiming the underlying inclusion-exclusion formula as new.

There is also a relevant recent development. Padhi's August 2026 preprint
reports holonomicity for every fixed $r,s$, while explicitly recording only a
partial result toward the specific A189281 operator~\cite{padhi}. Our arguments
neither assume that preprint's holonomicity theorem nor prove a particular
recurrence. A general existence theorem for a recurrence and a proof of a
specific guessed operator are distinct questions.
[Added 1 October 2026, batch 73O2: the intake confirmed that the arXiv record
of~\cite{padhi} exists under the title given in the bibliography, but did not
check its content. Nothing below depends on it.]

\subsection{The gap addressed here}

We supply an unconditional route from the counting problem to every coefficient
and a controlled remainder. The important point is not merely to fit more
terms. One must justify replacing an $n$-dependent inclusion-exclusion sum by an
infinite coefficient calculation, including the contributions of arbitrarily
many isolated selected edges. The exact stabilization theorem and a uniform
factorial-moment bound supply that justification.

The resulting contribution has four connected parts. First, the coefficients
attached to the two conjectural recurrences become theorems about the original
permutation counts. Second, the same proof gives every fixed pair of offsets
and the full violation distribution. Third, it explains why congruence classes
of $n$ do not enter any algebraic-order coefficient. Fourth, it gives an inverse
expansion with a precise interpretation for a discrete sequence.

The ProveIt repository was consulted as the project context. Its Fabius
asymptotic audit carefully separates a printed approximation, a necessary
correction, and an all-orders theorem~\cite{proveit}. Here the analogous
separation is between finite data, formal coefficient extraction, a uniform
asymptotic theorem, and an unproved stronger closed form. No declaration from
the repository is a mathematical dependency of the present proof, and no new
repository build is reported.

\subsection{A preview of the main theorem}

Put
\begin{equation}
 G^{(\theta)}_{r,s,n}(u)=\E(1+u)^{X^{(\theta)}_{r,s,n}}.
\end{equation}
The use of $1+u$, rather than $u$, makes the coefficients factorial moments.

\begin{theorem}[All-orders polynomial correction theorem]
\label{spf:thm:main}
For fixed positive integers $r,s$, $\theta\in\{1,2\}$, every integer $M\ge0$,
and every $0\le U<\infty$, there are polynomials
$B_J^{(\theta)}(u;r,s)\in\Q[u]$ such that, uniformly for $|u|\le U$,
\begin{equation}
 G^{(\theta)}_{r,s,n}(u)
 =e^{\theta u}\sum_{J=0}^{M}\frac{B_J^{(\theta)}(u;r,s)}{n^J}
   +\bigO_{M,r,s,U}(n^{-M-1}).
 \label{spf:eq:main}
\end{equation}
They are given by the finite formula \eqref{spf:eq:Bformula}. They satisfy
\begin{gather}
 B_0=1,\qquad \deg B_J\le2J,\qquad B_J(0)=0\quad(J\ge1),\\
 B_J(u;r,s)=B_J(u;s,r),\qquad (2J)!B_J(u;r,s)\in\Z[u],
 \label{spf:eq:arithmetic-preview}\\
 B_1^{(\theta)}(u;r,s)
 =\theta(1-r-s)u+\theta(1-\theta)u^2.
 \label{spf:eq:B1}
\end{gather}
In particular, for $c_J^{(\theta)}(r,s)=B_J^{(\theta)}(-1;r,s)$,
\begin{equation}
 \frac{A^{(\theta)}_{r,s}(n)}{n!}
 =e^{-\theta}\left(\sum_{J=0}^{M}\frac{c_J^{(\theta)}(r,s)}{n^J}
                   +\bigO(n^{-M-1})\right),
 \quad c_0=1,\quad c_1=\theta(r+s-\theta).
 \label{spf:eq:avoidance-expansion}
\end{equation}
\end{theorem}
The proof occupies \cref{spf:sec:tilings,spf:sec:stable,spf:sec:bounds,spf:sec:coefficients,spf:sec:proof}.
The arithmetic assertions are proved in \cref{spf:sec:arithmetic}.

\section{Selected violations and pairs of path tilings}
\label{spf:sec:tilings}

\subsection{Why paths are the right objects}

The directed graph on $\{1,\ldots,n\}$ with edges $i\to i+r$ is a disjoint union
of $r$ directed paths, when $n\ge r$. Their vertex counts are
$\lfloor n/r\rfloor$ and $\lceil n/r\rceil$. The value graph with edges
$j\to j+s$ is similarly a union of $s$ paths.

More generally, consider any directed path forest on the position vertices and
any directed path forest on the value vertices, both having $n$ vertices.
The oriented statistic counts position edges mapped to directed value edges by
a uniform bijection. The absolute statistic counts position edges mapped to
value edges in either direction. The proofs below apply to this more general
model as well.

Selecting a subset of the position edges gives a collection of nontrivial path
components and singleton vertices. For the oriented statistic, each selected
component must map in its prescribed order to a value-path interval. For the
absolute statistic there are exactly two choices of direction per nontrivial
component. A reversal of direction \emph{within} one component is impossible:
at the first change of sign it would repeat the preceding value. Thus the
orientation weight is $2^c$, not $2^k$, where $c$ is the number of components
and $k$ the number of selected edges.

\begin{definition}[Component profile]
A profile is a finitely supported sequence
$\mathbf b=(b_2,b_3,\ldots)$ of nonnegative integers. Define
\begin{equation}
 \begin{aligned}
 c&=\sum_{\ell\ge2}b_\ell,&
 k&=\sum_{\ell\ge2}(\ell-1)b_\ell,\\
 v&=\sum_{\ell\ge2}\ell b_\ell=k+c,&
 d&=k-c=\sum_{\ell\ge2}(\ell-2)b_\ell,
 \end{aligned}
 \label{spf:eq:profile}
\end{equation}
and $\bfac=\prod_{\ell\ge2}b_\ell!$.
The number $d$ is the \emph{defect}: dimers have defect zero, and each extra
edge in a component adds one unit of defect.
\end{definition}

For a position forest $\mathcal P$, let $C_{\mathcal P}(\mathbf b)$ be its
number of tilings with $b_\ell$ components of size $\ell\ge2$, and $n-v$
singletons. A tiling is simply a partition into consecutive intervals along
each path; unused edges may lie between neighboring intervals.

\begin{proposition}[Exact profile identity]
\label{spf:prop:exact}
For position and value forests $\mathcal P,\mathcal V$ on $n$ vertices,
\begin{equation}
 \E(1+u)^X
 =\sum_{\mathbf b:\,v\le n}
    \theta^c u^k C_{\mathcal P}(\mathbf b)C_{\mathcal V}(\mathbf b)
    \frac{(n-v)!\,\bfac}{n!}.
 \label{spf:eq:exact-profile}
\end{equation}
Infeasible profiles have zero tiling count.
\end{proposition}
\begin{proof}
For each permutation, expand $(1+u)^X$ as a sum over its selected violation
sets. Such a set determines its input tiling and its image tiling, with the
same component profile. Match the $b_\ell$ input components of each length to
the $b_\ell$ output components in $b_\ell!$ ways. Choose a direction for each
nontrivial component when $\theta=2$, and match the $n-v$ singletons
arbitrarily. This gives exactly
$\theta^c(n-v)!\bfac$ permutations for each pair of tilings, counted with the
selected set distinguished. Division by $n!$ proves the identity.

Additional violations outside the selected set are allowed. This is essential:
the identity counts marked subsets of violations, not permutations having
exactly that selected set.
\end{proof}

At $u=-1$, \eqref{spf:eq:exact-profile} is the matching-of-tilings
inclusion-exclusion formula of Spahn and Zeilberger~\cite{spahn}. Its extension
to general $u$ makes the later distributional statements possible.

\section{Exact stabilization on long path forests}
\label{spf:sec:stable}

Write $\fall{x}{j}=x(x-1)\cdots(x-j+1)$ and
$\rise{x}{j}=x(x+1)\cdots(x+j-1)$, with both empty products equal to one.
For a profile $\mathbf b$, put
\begin{equation}
 E_h(\mathbf b)=[z^h]\prod_{\ell\ge2}(1+(\ell-1)z)^{b_\ell}.
 \label{spf:eq:E}
\end{equation}
Define an integer polynomial
\begin{equation}
 F_r(n;\mathbf b)=
 \sum_{h=0}^{c}(-1)^h\rise{(r-1)}{h}E_h(\mathbf b)
                   \fall{(n-k-h)}{c-h}.
 \label{spf:eq:F}
\end{equation}
Its degree in $n$ is $c$ and its leading coefficient is one. For $r=1$, it
reduces to $\fall{(n-k)}c$.

\begin{theorem}[Stable tiling polynomial]
\label{spf:thm:stable}
Let $\mathcal P$ be a union of $r$ directed paths with vertex counts
$L_1,\ldots,L_r$ and total $n$. If each $L_i\ge k$, then
\begin{equation}
 C_{\mathcal P}(\mathbf b)=\frac{F_r(n;\mathbf b)}{\bfac}.
 \label{spf:eq:stable-C}
\end{equation}
Thus, in this range, the answer depends on $n,r,\mathbf b$, but not on the
individual path lengths.
\end{theorem}

\subsection{A formal coefficient identity}

The following residue calculation is included to make the stabilization
mechanism explicit. Only finitely many variables are needed for any specified
coefficient.

Let $z_2,z_3,\ldots$ be formal variables and let $W$, with zero constant term,
solve
\begin{equation}
 W=\sum_{\ell\ge2}z_\ell(1+W)^{-(\ell-1)}.
\end{equation}
Define
\begin{equation}
 \mathcal D=1+\sum_{\ell\ge2}(\ell-1)z_\ell(1+W)^{-\ell}.
\end{equation}
For every integer $N$ and every profile,
\begin{equation}
 [\mathbf z^{\mathbf b}]\frac{(1+W)^N}{\mathcal D}
     =\frac{\fall{(N-k)}c}{\bfac}.
 \label{spf:eq:lagrange-id}
\end{equation}
To prove this, replace every $z_\ell$ by $t z_\ell$ and put
$\phi(w)=\sum z_\ell(1+w)^{-(\ell-1)}$. Then $W=t\phi(W)$ and
$\mathcal D=1-t\phi'(W)$. Work first over the fraction field in finitely many $z_\ell$, so that
$\phi(0)$ is invertible; the resulting identity is polynomial coefficientwise.
The change of formal residue variable $t=w/\phi(w)$ gives
\begin{equation}
 [t^c]\frac{f(W)}{1-t\phi'(W)}=[w^c]f(w)\phi(w)^c.
\end{equation}
Indeed, the Jacobian factor cancels $1-t\phi'(W)$. The coefficient of
$\mathbf z^{\mathbf b}$ in $\phi(w)^c$ is
$c!(1+w)^{-k}/\bfac$. Taking $f(w)=(1+w)^N$ proves
\eqref{spf:eq:lagrange-id}.

\subsection{Proof of the stabilization theorem}

For a single path of length $L$, a profile with $v\le L$ has $L-k$ tiles in
all. Ordering those tiles gives
\begin{equation}
 C_L(\mathbf b)=\frac{(L-k)!}{(L-v)!\bfac}
              =\frac{\fall{(L-k)}c}{\bfac}.
 \label{spf:eq:onepath}
\end{equation}
If $k\le L<v$, the falling factorial is zero, again giving the correct count.
Consequently, through weighted edge degree $k\le L$, the actual tiling series
agrees with the formal surrogate $(1+W)^L/\mathcal D$.

For $r$ paths, all allocations of a total profile of edge degree $k$ to the
paths have local edge degree at most $k$. Under $L_i\ge k$, multiplying the
surrogates is therefore legitimate at the required coefficient. Their product
is
\begin{equation}
 \frac{(1+W)^n}{\mathcal D^r}.
 \label{spf:eq:forest-surrogate}
\end{equation}
Expand the additional factor $\mathcal D^{-(r-1)}$. A selected subprofile
$\mathbf a\le\mathbf b$, with $|\mathbf a|=h$, contributes
\begin{equation}
 \frac{(-1)^h\rise{(r-1)}h}{\mathbf a!}
 \prod_{\ell\ge2}(\ell-1)^{a_\ell}\mathbf z^{\mathbf a}
 \frac{(1+W)^{n-\sum\ell a_\ell}}{\mathcal D}.
\end{equation}
Applying \eqref{spf:eq:lagrange-id} to the remaining profile gives
\begin{equation}
 \frac{(-1)^h\rise{(r-1)}h
          \prod(\ell-1)^{a_\ell}}{\mathbf a!(\mathbf b-\mathbf a)!}
 \fall{(n-k-h)}{c-h}.
\end{equation}
Summing over subprofiles of size $h$ replaces the product of binomial
coefficients by $E_h(\mathbf b)/\bfac$. This is exactly
\eqref{spf:eq:F}--\eqref{spf:eq:stable-C}, proving \cref{spf:thm:stable}.
\hfill$\square$

\begin{corollary}[Stable factorial moments]
\label{spf:cor:stable-moments}
If the position forest has $r$ paths, the value forest has $s$ paths, every
path has at least $k$ vertices, and $2k\le n$, then
\begin{equation}
 \mu_k:=\E\binom Xk
 =\sum_{\mathbf b:\,\sum(\ell-1)b_\ell=k}
       \frac{\theta^c F_r(n;\mathbf b)F_s(n;\mathbf b)}
            {\bfac\,\fall nv}.
 \label{spf:eq:stable-moment}
\end{equation}
For the ordinary offset model, this is exact whenever
\begin{equation}
 k\le\min\left(\left\lfloor\frac nr\right\rfloor,
               \left\lfloor\frac ns\right\rfloor,
               \left\lfloor\frac n2\right\rfloor\right).
 \label{spf:eq:stable-range}
\end{equation}
\end{corollary}
\begin{proof}
Use \cref{spf:thm:stable} on both tiling counts in \eqref{spf:eq:exact-profile} and
extract $u^k$. The condition $2k\le n$ ensures $v=k+c\le n$ for every profile
in the displayed sum.
\end{proof}

As a first check, a single dimer has $F_r=n-r$, so, for
$n\ge\max(r,s,2)$,
\begin{equation}
 \E X=\frac{\theta(n-r)(n-s)}{n(n-1)}.
 \label{spf:eq:exact-mean}
\end{equation}
The same formula follows directly by counting the available input and output
edges.

\section{Uniform bounds and beyond-all-orders universality}
\label{spf:sec:bounds}

An exact formula for fixed moments is not, by itself, enough to justify an
asymptotic expansion of a probability. A bound uniform in the moment order is
needed.

\begin{lemma}[Uniform factorial-moment bound]
\label{spf:lem:moment-bound}
For either statistic and arbitrary path forests on $n\ge1$ vertices,
\begin{equation}
 0\le\mu_k\le\frac{(\theta e^2)^k}{k!}\qquad(k\ge0).
 \label{spf:eq:moment-bound}
\end{equation}
Consequently, for $|u|\le U$ and $a=\theta e^2U$,
\begin{equation}
 \left|G(u)-\sum_{k=0}^{K}\mu_k u^k\right|
 \le e^a\frac{a^{K+1}}{(K+1)!}.
 \label{spf:eq:moment-tail}
\end{equation}
\end{lemma}
\begin{proof}
Fix a selected set of $k$ position edges, with $c$ components and $v=k+c$
vertices. There are at most $\theta^c n^c$ ways to assign starts and directions
to its components in the value forest, even before excluding overlaps. Hence
the probability that all the selected edges are violations is at most
$\theta^c n^c/\fall nv$.

For $0\le v\le n$, the geometric mean of the largest $v$ factors in $n!$ is at
least the geometric mean of all $n$ factors. Also
$n!\ge(n/e)^n$. Thus $\fall nv\ge(n/e)^v$, and the last probability is at most
$\theta^k e^{2k}n^{-k}$, since $c\le k$ and $v\le2k$.
There are at most $n^k/k!$ possible selected edge sets. Summing proves
\eqref{spf:eq:moment-bound}. Summing its exponential-series tail proves
\eqref{spf:eq:moment-tail}. For $k=0$ the assertion is immediate.
\end{proof}

\begin{theorem}[Quantitative independence of long path lengths]
\label{spf:thm:universality}
Consider two position--value forest pairs, each on $n$ vertices, with the same
numbers $r$ and $s$ of paths. Suppose every path in both pairs has at least $K$
vertices and $2K\le n$. Let $X$ and $\widetilde X$ be the corresponding
statistics for the same $\theta$. For every $R>1$, with
$a=\theta e^2(R+1)$,
\begin{equation}
 \TV(\mathcal L(X),\mathcal L(\widetilde X))
 \le \frac{R}{R-1}\,e^a\frac{a^{K+1}}{(K+1)!}.
 \label{spf:eq:TV-universality}
\end{equation}
If all path lengths are at least $\eta n$ for a fixed $\eta>0$, the right-hand
side can be chosen to be
$\exp(-\eta' n\log n+\bigO(n))$ for every fixed
$0<\eta'<\min(\eta,1/2)$.
\end{theorem}
\begin{proof}
By Corollary~\ref{spf:cor:stable-moments}, the two sets of factorial moments agree through
order $K$. On the circle $|t|=R$, use $u=t-1$ in
\eqref{spf:eq:moment-tail}; the two ordinary probability generating functions differ
by at most $2e^a a^{K+1}/(K+1)!$. Cauchy's coefficient inequality bounds the
difference of the probabilities at $j$ by this quantity times $R^{-j}$.
Sum over $j\ge0$ and divide by two to obtain \eqref{spf:eq:TV-universality}.
Finally, take $K=\lfloor\eta'n\rfloor$ and use
$\log(K!)=K\log K-K+\bigO(\log K)$.
\end{proof}

This statement is stronger than agreement of a few asymptotic coefficients.
It concerns the entire probability law. It also clarifies a limitation: the
bound does not identify the leading exponentially small difference between
two particular forest geometries. That finer sector is left open.

\section{A finite formula for every asymptotic coefficient}
\label{spf:sec:coefficients}

\subsection{Normalization and the defect filtration}

Put $x=1/n$ and define
\begin{align}
 f_{r,\mathbf b}(x)
 &=\sum_{h=0}^{c}(-1)^h\rise{(r-1)}h E_h(\mathbf b)x^h
       \prod_{a=0}^{c-h-1}(1-(k+h+a)x),\\
 D_v(x)&=\prod_{a=0}^{v-1}(1-ax),\\
 Q_{\mathbf b}(x)&=\frac{f_{r,\mathbf b}(x)f_{s,\mathbf b}(x)}{D_v(x)}.
 \label{spf:eq:Q}
\end{align}
Then $f_{r,\mathbf b}(1/n)=n^{-c}F_r(n;\mathbf b)$, and a stabilized profile
contributes
\begin{equation}
 \frac{\theta^c u^k}{\bfac}\,n^{-d}Q_{\mathbf b}(1/n)
 \label{spf:eq:normalized-profile}
\end{equation}
to the generating function. This displays why defect $d>J$ cannot contribute
to the coefficient of $n^{-J}$.

Separate the dimers from the other components:
\begin{equation}
 m=b_2,\qquad \boldsymbol\beta=(b_3,b_4,\ldots).
\end{equation}
For the latter profile write
\begin{equation}
 C=\sum_{\ell\ge3}\beta_\ell,\quad
 K=\sum_{\ell\ge3}(\ell-1)\beta_\ell,\quad
 d=K-C,\quad \betafac=\prod_{\ell\ge3}\beta_\ell!.
\end{equation}
Thus $c=m+C$, $k=m+K$, $v=2m+K+C$.
The profiles $\boldsymbol\beta$ of a given defect $d$ are in bijection with
ordinary integer partitions of $d$: a component of length $\ell$ contributes a
part $\ell-2$.

\begin{lemma}[Polynomial dependence on the number of dimers]
\label{spf:lem:polynomial}
For fixed $\boldsymbol\beta,j,r,s$,
\begin{equation}
 P_{\boldsymbol\beta,j}(m):=[x^j]Q_{(m,\boldsymbol\beta)}(x)
 \label{spf:eq:P}
\end{equation}
is a polynomial in $m$ over $\Q$ of degree at most $2j$. For positive integer
$r,s$, it takes integer values at every $m\ge0$.
\end{lemma}
\begin{proof}
The coefficient $E_h(m,\boldsymbol\beta)$ is a polynomial in $m$ of degree at
most $h$, since the dimer contribution is $(1+z)^m$.
For a product of consecutive linear factors
\begin{equation}
 \prod_{a=0}^{m+C-h-1}(1-(m+K+h+a)x),
\end{equation}
the coefficient of $x^q$ is polynomial in $m$ of degree at most $2q$.
One way to see the bound is to take its formal logarithm. The coefficient of
$x^i$ is a constant multiple of a sum of $i$th powers over a consecutive
interval with affine endpoints, and therefore has degree at most $i+1$.
Exponentiating shows degree at most $2q$, because a product of power sums of
total $x$-degree $q$ has degree at most $q+q$.
The same argument applies to $D_v(x)^{-1}$.

In a term carrying the additional factor $E_hx^h$, the degree at coefficient
$x^j$ is at most $h+2(j-h)\le2j$. Products preserve the total bound. For the
few small nonnegative $m$ for which $c<h$, $E_h=0$; thus the polynomial
extension gives the same zero contribution. Equivalently one can define the
consecutive power sums by their polynomial extensions before multiplying by
$E_h$.

Finally, for nonnegative $m$, the numerator polynomials in \eqref{spf:eq:Q} have
integer coefficients and $D_v^{-1}\in\Z[[x]]$. Hence every value in
\eqref{spf:eq:P} is an integer.
\end{proof}

\subsection{Poisson summation by finite differences}

Let $\Delta P(m)=P(m+1)-P(m)$. For a polynomial of degree at most $q$,
\begin{equation}
 P(m)=\sum_{t=0}^{q}\frac{\Delta^tP(0)}{t!}\fall mt.
\end{equation}
Therefore the elementary identity
\begin{equation}
 \sum_{m\ge0}\frac{z^m}{m!}P(m)
  =e^z\sum_{t=0}^{q}\frac{\Delta^tP(0)}{t!}z^t
 \label{spf:eq:poisson-transform}
\end{equation}
turns an infinite dimer sum into a finite polynomial. It follows either by
summing each falling factorial, or by differentiating $e^z$ exactly $t$ times.

\begin{definition}[Coefficient formula]
\label{spf:def:B}
For $J\ge0$, define
\begin{equation}
 \boxed{\displaystyle
 B_J^{(\theta)}(u;r,s)=
 \sum_{d=0}^{J}
 \ \sum_{\substack{\boldsymbol\beta\ge0\\
                   \sum_{\ell\ge3}(\ell-2)\beta_\ell=d}}
 \frac{\theta^C u^K}{\betafac}
 \sum_{t=0}^{2(J-d)}
 \frac{\Delta^t P_{\boldsymbol\beta,J-d}(0)}{t!}(\theta u)^t .}
 \label{spf:eq:Bformula}
\end{equation}
All sums on the right are finite.
\end{definition}

Only evaluations at $m=0,1,\ldots,2(J-d)$ are required. Their sufficiency is a
proved degree bound, not a polynomial guess. No term $A(n)$ and no recurrence
for $A(n)$ is needed to compute $B_J$.

\section{Proof of the uniform all-orders expansion}
\label{spf:sec:proof}

We now justify the asymptotic interpretation of \eqref{spf:eq:Bformula}. This is
the analytic step that a purely formal expansion of the profile formula would
leave unresolved.

\begin{lemma}[Uniform analytic bound for a profile]
\label{spf:lem:Qbound}
For every profile of edge degree $k$, put
$\rho_k=1/[8(k+1)^2]$. There is a constant $C_{r,s}$, independent of the
profile and $k$, such that
\begin{equation}
 \sup_{|x|\le\rho_k}|Q_{\mathbf b}(x)|\le C_{r,s}.
 \label{spf:eq:Qbound}
\end{equation}
One may take $C_{r,s}=e\,2^{r+s}$. In particular,
\begin{equation}
 |[x^j]Q_{\mathbf b}(x)|\le C_{r,s}[8(k+1)^2]^j,
 \label{spf:eq:Qcoeffbound}
\end{equation}
and, if $n\ge16(k+1)^2$,
\begin{equation}
 \left|Q_{\mathbf b}(1/n)-\sum_{j=0}^{q}[x^j]Q_{\mathbf b}(x)n^{-j}\right|
 \le2C_{r,s}\left(\frac{8(k+1)^2}{n}\right)^{q+1}.
 \label{spf:eq:Qremainder}
\end{equation}
\end{lemma}
\begin{proof}
The elementary symmetric function bound
$E_h(\mathbf b)\le k^h/h!$ follows by expanding the $h$th power of the sum
of the component edge lengths. Every integer appearing in a numerator factor
is at most $2k$, and there are at most $k$ such factors. Thus, on the stated
disk,
\begin{equation}
 |f_{r,\mathbf b}(x)|
 \le e^{2k^2\rho_k}
       \sum_{h\ge0}\frac{\rise{(r-1)}h}{h!}(k\rho_k)^h
 =e^{2k^2\rho_k}(1-k\rho_k)^{-(r-1)}.
\end{equation}
For $r=1$, the last sum is simply one. Also $v\le2k$ and $a\rho_k<1/2$ for
$0\le a<v$, so
\begin{equation}
 |D_v(x)^{-1}|\le\exp\left(2\rho_k\sum_{a=0}^{v-1}a\right)
                 \le e^{4k^2\rho_k}.
\end{equation}
Combining these bounds gives \eqref{spf:eq:Qbound} with the stated generous
constant. The denominator has no zero on the disk. Cauchy's coefficient bound
then proves \eqref{spf:eq:Qcoeffbound}, and its geometric tail at
$1/n\le\rho_k/2$ proves \eqref{spf:eq:Qremainder}.
\end{proof}

\begin{proof}[Proof of \cref{spf:thm:main}: analytic assertion]
Fix $M,r,s,U$, and for sufficiently large $n$ choose
\begin{equation}
 T_n=\left\lfloor\frac{(2M+4)\log n}{\log\log n}\right\rfloor.
\end{equation}
The factorial-moment tail \eqref{spf:eq:moment-tail} with $K=T_n$ is
$n^{-(2M+4)+o(1)}$, hence is $\bigO(n^{-M-1})$. For fixed offsets, every
position and value path has more than $T_n$ vertices eventually. Also
$2T_n\le n$ and $16(T_n+1)^2\le n$. Thus all retained moments are given
exactly by \eqref{spf:eq:stable-moment}, and Lemma~\ref{spf:lem:Qbound} applies to every
retained profile.

\emph{High defect.} The sum of absolute profile weights, marked by their defect,
has the convergent generating function
\begin{equation}
 \sum_{\mathbf b}\frac{\theta^c U^k}{\bfac}z^d
 =\exp\left(\sum_{q\ge1}\theta U^q z^{q-1}\right)
 =\exp\left(\frac{\theta U}{1-Uz}\right).
 \label{spf:eq:defect-majorant}
\end{equation}
For $U=0$ the assertion is immediate; otherwise use a disk strictly inside
$|z|<1/U$. Its Taylor tail of degree $M$ at $z=1/n$ is
$\bigO(n^{-M-1})$. Together with \eqref{spf:eq:Qbound}, this bounds all retained
profiles of defect $d\ge M+1$ by $\bigO(n^{-M-1})$.

\emph{Low defect.} For each $d\le M$, there are only finitely many
$\boldsymbol\beta$. Expand $Q_{(m,\boldsymbol\beta)}$ through order $M-d$.
By \eqref{spf:eq:Qremainder}, after multiplying by $n^{-d}$ the error is bounded
by $n^{-M-1}$ times a fixed polynomial in $m$ times
$(\theta U)^m/m!$, with a constant depending only on the fixed parameters
and $\boldsymbol\beta$. The sum over $m$ is finite uniformly in $n$. This is
why no extra power of $\log n$ appears in the remainder.

We may now extend the retained dimer sums to $m=\infty$. Indeed, for each
retained coefficient, \eqref{spf:eq:Qcoeffbound} bounds the omitted summands by a
polynomial in $m$ times $(\theta U)^m/m!$. Since $k=m+K>T_n$ implies
$m>T_n-K$, their tails are again
$n^{-(2M+4)+o(1)}$. Summing each polynomial in $m$ by
\eqref{spf:eq:poisson-transform} gives exactly \eqref{spf:eq:Bformula}. The total error
is $\bigO(n^{-M-1})$, uniformly on $|u|\le U$. This proves
\eqref{spf:eq:main} and, on setting $u=-1$, \eqref{spf:eq:avoidance-expansion}.
\end{proof}

The proof also works for any fixed numbers $r,s$ of paths whose shortest
length is bounded below by a positive multiple of $n$. It does not assert
uniformity for arbitrarily growing offsets or for forests with a bounded
short path.
[Added 1 October 2026, batch 74: Part~II proves the expansion uniformly
over forest shapes, on the disk $|1+u|\le1$, whenever the shortest path has
at least $K_n\asymp\log n/\log\log n$ vertices
(\cref{fga:thm:allorders,fga:cor:shortest}); \cref{fga:rem:partI-shapes}
notes that the proof above already gives this on every disk $|u|\le U$
with $T_n$ in place of $K_n$.]

\section{Arithmetic, first corrections, and the full distribution}
\label{spf:sec:arithmetic}

\subsection{Degree and denominator control}

\begin{proposition}[Arithmetic of the correction polynomials]
\label{spf:prop:arithmetic}
The structural assertions \eqref{spf:eq:arithmetic-preview} hold.
\end{proposition}
\begin{proof}
For a defect-$d$ profile, $C\le d$ and $K\le2d$. A term in
\eqref{spf:eq:Bformula} has $u$-degree at most
$K+2(J-d)\le2J$. The formula is symmetric in $r,s$ because $Q_{\mathbf b}$ is.
For $J=0$ it gives $B_0=1$. When $J>0$, any nonempty
$\boldsymbol\beta$ contributes positive $u$-degree. For the empty profile the
constant term is $P_{\varnothing,J}(0)=0$, since the completely empty
$\mathbf b$ has $Q_{\mathbf b}=1$. Thus $B_J(0)=0$.

By Lemma~\ref{spf:lem:polynomial}, every $\Delta^tP(0)$ is an integer. A summand's only
possible denominator is $t!\betafac$. This divides $(t+C)!$, by the
multinomial coefficient, and
\begin{equation}
 t+C\le2(J-d)+d\le2J.
\end{equation}
Hence $t!\betafac$ divides $(2J)!$. Multiplying the finite sum by $(2J)!$
therefore gives a polynomial over $\Z$.
\end{proof}

This denominator bound is unconditional. It is deliberately weaker than
claiming all the oriented coefficients themselves are integers; that stronger
phenomenon is discussed in \cref{spf:sec:collapse}.

\subsection{The first two corrections}

For the dimer-only profile,
\begin{equation}
 P_{\varnothing,1}(m)=-m^2+(2-r-s)m
 =-\fall m2+(1-r-s)m.
\end{equation}
The only defect-one profile is a single component of length three, with
$C=1$, $K=2$, and $P_{\boldsymbol\beta,0}=1$. Applying
\eqref{spf:eq:poisson-transform} immediately gives \eqref{spf:eq:B1}. For $\theta=1$,
a calculation of the four profile types of defect at most two gives
\begin{equation}
 B_2^{(1)}(u;r,s)
 =(r-1)(s-1)u+\binom{r+s-1}{2}u^2.
 \label{spf:eq:B2-oriented}
\end{equation}
For completeness, the required dimer-only polynomial is
\begin{equation}
 \begin{split}
 P_{\varnothing,2}(m)
 ={}&\frac12m^4+(r+s-3)m^3\\
 &+\frac{r^2+2rs+s^2-5r-5s+7}{2}m^2
   +\frac{-r^2+r-s^2+s}{2}m.
 \end{split}
\end{equation}
For one length-three component,
$P_{\boldsymbol\beta,1}(m)=-m^2-(r+s)m+3-2r-2s$.
The defect-two profiles are one length-four component or two length-three
components; both have $P_{\boldsymbol\beta,0}=1$. These data and
\eqref{spf:eq:Bformula} verify \eqref{spf:eq:B2-oriented} directly.

In the absolute case a useful general expression is
\begin{equation}
 \begin{split}
 B_2^{(2)}(u;r,s)={}&2(r-1)(s-1)u
 +2\bigl((r+s)^2-r-s-1\bigr)u^2\\
 &+2(2r+2s-1)u^3+2u^4.
 \label{spf:eq:B2-absolute}
 \end{split}
\end{equation}
Thus \eqref{spf:eq:tauraso} is recovered as a special case of the first correction,
while the theorem supplies every later order.

\begin{remark}[Offsets $(1,1)$: two classical sequences]
\label{spf:rem:offsets-one-one}
[Added 1 October 2026, batch 73O2.] The manuscript does not discuss the
smallest offsets, which give two classical sequences. The intake ran
the shipped coefficient generator \texttt{code/path\_forests.py} on a copy with
\texttt{--r 1 --s 1 --order 10}.
\begin{itemize}
\item \emph{Oriented case, no successions.} Here
 $A^{(1)}_{1,1}(n)$ counts permutations of $[n]$ with no $\pi(i+1)=\pi(i)+1$,
 which is $\textnormal{A000255}(n-1)=D_n+D_{n-1}$ with $D_n$ the derangement
 numbers~\cite{oeis000255}. The formula gives
 $c_J^{(1)}(1,1)=1,1,0,0,\ldots,0$ for $J=0,\ldots,10$. This is consistent
 with $(D_n+D_{n-1})/n!=e^{-1}(1+1/n)+\bigO(1/n!)$, so every later
 coefficient is zero as well.
\item \emph{Absolute case, Hertzsprung's problem.} Here
 $A^{(2)}_{1,1}(n)$ is \textnormal{A002464}, the number of permutations with
 no rising or falling succession~\cite{oeis002464}. The formula gives
 \begin{gather*}
 c_J^{(2)}(1,1)=1,\ 0,\ -2,\ -\tfrac{10}{3},\ -6,\ -\tfrac{154}{15},\
 -\tfrac{88}{9},\\
 \tfrac{5336}{105},\ \tfrac{1612}{3},\
 \tfrac{2098234}{567},\ \tfrac{36500686}{1575}\qquad(J=0,\ldots,10).
 \end{gather*}
 These are, term for term, the Abramson--Moser expansion displayed in the
 OEIS entry (computed by Kotesovec, 2011, extended 2020), through $n^{-10}$.
\end{itemize}
Both cases agree with $c_1=\theta(r+s-\theta)$ of \eqref{spf:eq:avoidance-expansion}.
As an independent check, the intake computed both sequences exactly to
$n=800$, from the recurrences printed in the two OEIS entries (themselves
checked against brute-force enumeration for $n\le8$). It found the scaled
residuals $n^{M+1}\bigl(e^{\theta}A(n)/n!-\sum_{j\le M}c_jn^{-j}\bigr)$ for
$M=4,8$ approaching $c_{M+1}$ at $n=200,400,800$. These finite checks are
evidence; the proof is \cref{spf:thm:main}. The theorem covers the pair
$(1,1)$, and the manuscript's brute-force tests already include it (it is one
of the six offset pairs of \cref{spf:sec:computation}). Only the comparison
with the two OEIS entries is new here.
\end{remark}

\subsection{Poisson corrections for all probabilities}

The zero-violation event is only one coefficient of the probability generating
function. Write $B_J(u)=\sum_h b_{J,h}u^h$ and define
\begin{equation}
 H_h(q;\theta)=\sum_{i=0}^{\min(h,q)}
      \binom hi(-1)^{h-i}\frac{\fall qi}{\theta^i}.
 \label{spf:eq:Charlier}
\end{equation}
These are a normalization of the Charlier correction polynomials.
Direct coefficient extraction gives
\begin{equation}
 [t^q](t-1)^h e^{\theta(t-1)}
 =e^{-\theta}\frac{\theta^q}{q!}H_h(q;\theta).
\end{equation}
Consequently define the signed approximation
\begin{equation}
 p_M(q)=e^{-\theta}\frac{\theta^q}{q!}
   \sum_{J=0}^{M}\frac1{n^J}\sum_h b_{J,h}H_h(q;\theta).
 \label{spf:eq:pmf}
\end{equation}

\begin{corollary}[An all-orders approximation of the entire law]
\label{spf:cor:distribution}
For fixed $M,r,s,\theta$,
\begin{equation}
 \sum_{q\ge0}\left|\Pp(X^{(\theta)}_{r,s,n}=q)-p_M(q)\right|
 =\bigO(n^{-M-1}).
 \label{spf:eq:l1}
\end{equation}
In particular the law converges in total variation to
$\mathrm{Poisson}(\theta)$, with error $\bigO(n^{-1})$.
\end{corollary}
\begin{proof}
Apply \cref{spf:thm:main} on $|t|=R>1$, using $u=t-1$. Cauchy's inequality gives
an error $\bigO(n^{-M-1}R^{-q})$ in the coefficient of $t^q$. Summing the
geometric series proves \eqref{spf:eq:l1}.
\end{proof}

For $M\ge1$, $p_M$ need not be nonnegative at every $q$, so it is a signed
asymptotic approximation, not a new probability model. It nevertheless has
total mass one, because $B_0(0)=1$ and $B_J(0)=0$ for $J>0$.

\begin{remark}[A clique-model analogue]
\label{spf:rem:clique-analogue}
[Added 1 October 2026, batch 73O2, reciprocal note.] The companion report
\texttt{a330266-balanced-smirnov-poisson} of this collection (batch~73,
manuscript~57) studies the number $X_{n,k}$ of equal-rank adjacencies in a
uniform shuffle of $kn$ cards, $k$ in each of $n$ ranks. There the value
relation is $n$ cliques $K_k$, not a fixed number of long paths, so the
hypotheses of this report do not apply, and the limit is
$\mathrm{Poisson}(\lambda)$ with $\lambda=k-1$. Its first-order local law
(Corollary~8.3 there, label \texttt{bsw:cor:local}),
\begin{equation*}
 \Pp(X_{n,k}=j)=e^{-\lambda}\frac{\lambda^j}{j!}
 \left[1-\frac{(\lambda-j)^2-j}{2kn}+\bigO_{k,j}(n^{-2})\right],
\end{equation*}
is the analogue of the case $M=1$ of \eqref{spf:eq:pmf}--\eqref{spf:eq:l1}.
That report proves it at first order only, and the analytic step behind its
expansion is sketched; its Remark~6.3 names a uniform factorial-moment bound
of the type of Lemma~\ref{spf:lem:moment-bound} as the device that would make
it rigorous.
\end{remark}

\section{The two OEIS expansions, unconditionally}
\label{spf:sec:OEIS-results}

\subsection{A189281: oriented distance-two avoidance}

The first correction polynomials for $r=s=2$, $\theta=1$ are
\begin{align}
 B_0&=1,& B_1&=-3u,& B_2&=u+3u^2,\\
 B_3&=u+u^2-u^3,&
 B_4&=u-u^2-2u^3,&
 B_5&=u+u^2-3u^3,\\
 B_6&=u+23u^2-4u^3,&
 B_7&=u+121u^2+19u^3.
\end{align}
Further examples are
\begin{align}
 B_8&=u+479u^2+378u^3+24u^4,\\
 B_9&=u+1681u^2+3737u^3+672u^4+24u^5,\\
 B_{10}&=u+5543u^2+28936u^3+11088u^4+984u^5+24u^6.
\end{align}
Setting $u=-1$ in these polynomials proves
\begin{equation}
 \begin{split}
 \frac{\Aone(n)}{n!}=e^{-1}\biggl(&1+\frac3n+\frac2{n^2}+\frac1{n^3}
 +\frac3{n^5}+\frac{26}{n^6}+\frac{101}{n^7}+\frac{124}{n^8}\\
 &-\frac{1409}{n^9}-\frac{13266}{n^{10}}+\bigO(n^{-11})\biggr).
 \end{split}
 \label{spf:eq:A189281}
\end{equation}
The coefficient of $n^{-4}$ is exactly zero. The terms through $n^{-10}$
agree with the recurrence-dependent list in the OEIS~\cite{oeis189281};
the proof above does not assume that recurrence.

\subsection{A110128: absolute distance-two avoidance}

For $r=s=2$, $\theta=2$, the first polynomials are
\begin{align}
 B_1&=-6u-2u^2,\\
 B_2&=2u+22u^2+14u^3+2u^4,\\
 B_3&=2u+6u^2-60u^3-56u^4-16u^5-\frac43u^6,\\
 B_4&=2u-18u^2-122u^3+74u^4+154u^5+70u^6+12u^7+\frac23u^8.
\end{align}
It follows that
\begin{equation}
 \begin{split}
 \frac{\Atwo(n)}{n!}=e^{-2}\biggl(&1+\frac4n+\frac8{n^2}
  +\frac{68}{3n^3}+\frac{242}{3n^4}+\frac{1692}{5n^5}\\
 &+\frac{72802}{45n^6}+\frac{2725708}{315n^7}
  +\frac{16083826}{315n^8}\\
 &+\frac{186091480}{567n^9}
  +\frac{32213578294}{14175n^{10}}
  +\bigO(n^{-11})\biggr).
 \end{split}
 \label{spf:eq:A110128}
\end{equation}
This recovers the higher list attributed to the guessed operator in the
OEIS~\cite{oeis110128}, again by a direct argument from the permutation model.
The fractions here also show why an integer-coefficient assertion cannot be
transferred indiscriminately from the oriented case to the absolute case.

\subsection{Additional coefficients and how to use them}

\begin{table}[htbp]
\centering
\small
\renewcommand{\arraystretch}{1.22}
\begin{tabular}{r r r}
\toprule
$J$ & $c_J^{(1)}(2,2)$ & $c_J^{(2)}(2,2)$\\
\midrule
11 & $-59103$ & $2620164054016/155925$\\
12 & $-9448$ & $61699558213516/467775$\\
13 & $2525459$ & $6658719787519636/6081075$\\
14 & $26632378$ & $136086447719175784/14189175$\\
15 & $152270749$ & $5133067332126139204/58046625$\\
16 & $106221948$ & $546416517735316663856/638512875$\\
\bottomrule
\end{tabular}
\caption{Exact evaluations of the finite coefficient formula beyond the
order-ten lists. The accompanying JSON includes every polynomial $B_J$ through
order sixteen, not only its value at $u=-1$.}
\label{spf:tab:extra}
\end{table}

For any fixed $M$, adding terms through $c_Mn^{-M}$ gives a remainder
$\bigO(n^{-M-1})$ inside the parentheses. This is a Poincar\'e statement with
$M$ fixed as $n\to\infty$. It does not establish convergence of the infinite
series, justify setting $M$ proportional to $n$, or identify an optimal
truncation rule. The exact finite coefficients in \cref{spf:tab:extra} are supplied
with their generating code rather than inferred from a numerical fit.

\section{Lambert-\texorpdfstring{$W$}{W} inversion with a discrete interpretation}
\label{spf:sec:inverse}

\begin{remark}[Relation to the repository's inversion apparatus]
\label{spf:rem:inverse-apparatus}
[Added 1 October 2026, batch 73O2.] This section re-derives, for its own
forward map, machinery that the canonical volume \emph{Transseries and
Inversion} of the ProveIt repository already contains~\cite{tai}, and the
manuscript does not cite it. The forward map
$e^{-\theta}\Gamma(n+1)C(1/n)$ is the gamma carrier of that volume's chapter
on the increasing inverse of Gamma (theorem ``All-orders expansion of
$\Gamma_+^{-1}$'', label \texttt{p6:thm:gamma}), perturbed by a power series
in $1/n$, which is the situation of its theorem ``Perturbed inversion''
(\texttt{p0:thm:perturbed-inversion}). The volume's chapter ``The
subfactorial'' (\texttt{p8:sec:top}) treats the closest analogue, the
derangement numbers, by citing that carrier. Normalizations differ: here
$z=L/W_0(L/e)$ with $L=\log y$, and the constant $\tfrac12\log(2\pi)$ sits in
$\kappa_\theta$; the volume uses $R=\log y-\tfrac12\log(2\pi)$ and
$T=R/W_0(R/e)$. The discussion of integer thresholds at the end of the
section is an instance of the volume's staircase theorem
(\texttt{p0:thm:staircase}). \Cref{spf:thm:inverse-first,spf:thm:inverse-all}
are therefore printed as a self-contained second route for this family of
sequences. No novelty is claimed for the inversion method; what is specific
to this report is the forward input, the coefficients $c_j$ of
\cref{spf:thm:main}.
\end{remark}

\subsection{Logarithmic coefficients}

For either sequence, or any fixed $r,s,\theta$, suppress the parameters and
write
\begin{equation}
 \frac{A(n)}{n!}\sim e^{-\theta}C(1/n),\qquad
 C(x)=1+\sum_{j\ge1}c_jx^j.
\end{equation}
The right side is a formal asymptotic series. Define
$\log C(x)=\sum_{j\ge1}\ell_jx^j$. An explicit coefficient recursion is
\begin{equation}
 \ell_j=c_j-\frac1j\sum_{k=1}^{j-1}k\ell_k c_{j-k}.
 \label{spf:eq:log-recursion}
\end{equation}
Combining the proved probability expansion with Stirling's expansion, in its
usual remainder-controlled form~\cite{dlmfGamma}, yields
\begin{equation}
 \log A(n)
 =n\log n-n+\tfrac12\log n+\kappa_\theta
       +\sum_{j=1}^{M}a_jn^{-j}+\bigO(n^{-M-1}),
 \label{spf:eq:logA}
\end{equation}
where
\begin{equation}
 \kappa_\theta=\tfrac12\log(2\pi)-\theta,\qquad
 a_j=\ell_j+\frac{\mathsf B_{j+1}}{j(j+1)}.
 \label{spf:eq:a-coeff}
\end{equation}
Here $\mathsf B_m$ are Bernoulli numbers; the additional term vanishes for
even $j$. For A189281,
\begin{equation}
 \ell_1=3,\quad\ell_2=-\tfrac52,\quad\ell_3=4,\quad\ell_4=-\tfrac{29}{4},
 \qquad a_1=\tfrac{37}{12},\ a_2=-\tfrac52,\ a_3=\tfrac{1439}{360}.
\end{equation}
For A110128, $a_1=49/12$.

\subsection{The first explicit inverse terms}

Given $y=A(n)$, let $L=\log y$ and
\begin{equation}
 z=\frac{L}{W_0(L/e)},\qquad w=\log z.
 \label{spf:eq:Lambert}
\end{equation}
The principal Lambert function is characterized by $W_0(x)e^{W_0(x)}=x$
for $x>0$~\cite{dlmfW}. Thus $z(\log z-1)=L$ exactly.

\begin{theorem}[First inverse corrections]
\label{spf:thm:inverse-first}
Along the sequence values $y=A(n)$ as $n\to\infty$,
\begin{equation}
 \boxed{\displaystyle
 n=z-\frac12-\frac{\kappa_\theta}{w}
     -\frac1z\left(\frac{c_1-1/24}{w}
                    +\frac{\kappa_\theta^2}{2w^3}\right)
     +\bigO\left(\frac1{z^2w}\right).}
 \label{spf:eq:inverse-first}
\end{equation}
For A189281 the coefficient $c_1-1/24$ is $71/24$; for A110128 it is $95/24$.
\end{theorem}
\begin{proof}
Set $n=z+D$. The derivative of $x\log x-x$ at $z$ is $w$, and its second
derivative is $1/z$. The constant-order terms of \eqref{spf:eq:logA}, after
subtracting $L$, give
$wD+\tfrac12w+\kappa_\theta=0$, hence
$D=d_0=-1/2-\kappa_\theta/w$ at that order. The coefficient of $z^{-1}$ is
\begin{equation}
 w d_1+\frac{d_0^2+d_0}{2}+a_1.
\end{equation}
Since $a_1=c_1+1/12$, this gives
\begin{equation}
 d_1=-\frac{c_1-1/24}{w}-\frac{\kappa_\theta^2}{2w^3}.
\end{equation}
The residual in the logarithmic equation is $\bigO(z^{-2})$; division by a
derivative asymptotic to $w$ gives the asserted index error. A precise
all-orders version of this argument follows next.
\end{proof}

\subsection{An all-orders inverse recursion}

Put $t=1/z$ and seek $D(t,w)=\sum_{j\ge0}d_j(w)t^j$. Define
\begin{equation}
 \begin{split}
 \mathcal H(t,D)={}&
 \frac{(1+tD)\log(1+tD)-tD}{t}
 +\frac12\log(1+tD)\\
 &+\sum_{j\ge1}a_jt^j(1+tD)^{-j}.
 \end{split}
 \label{spf:eq:H-inverse}
\end{equation}
The quotient by $t$ is a formal series divisible by $t$, so no singular
constant term occurs. The inverse equation is
\begin{equation}
 wD+\tfrac12w+\kappa_\theta+\mathcal H(t,D)=0.
\end{equation}
It gives the triangular recursion
\begin{align}
 d_0(w)&=-\tfrac12-\kappa_\theta/w,\\
 d_j(w)&=-\frac1w[t^j]\mathcal H\left(t,\sum_{q=0}^{j-1}d_q(w)t^q\right)
             \qquad(j\ge1).
 \label{spf:eq:inverse-recursion}
\end{align}
Only $a_1,\ldots,a_j$ enter the $j$th step.

\begin{theorem}[All-orders index inversion]
\label{spf:thm:inverse-all}
For every fixed $J\ge0$, the functions in \eqref{spf:eq:inverse-recursion} are
rational in $w$, with coefficients in $\Q[\kappa_\theta]$ for the rational
forward coefficients considered here, and
\begin{equation}
 n=z+\sum_{j=0}^{J}\frac{d_j(w)}{z^j}
       +\bigO\left(\frac1{z^{J+1}w}\right),
 \qquad y=A(n).
 \label{spf:eq:inverse-all}
\end{equation}
Moreover, $d_0=\bigO(1)$ and $d_j=\bigO(1/w)$ for $j\ge1$ as $w\to\infty$.
\end{theorem}
\begin{proof}
The recursion first proves the algebraic form of the coefficients by
induction. Since $\mathcal H$ has no $t^0$ term, its coefficient of $t^j$
depends only on the earlier $d_q$, and stays bounded when those earlier
coefficients are bounded. Dividing by $w$ proves the stated size bounds.

Fix $J$ and retain sufficiently many terms of \eqref{spf:eq:logA}. Substitution
of the truncated $D$ cancels every term through $z^{-J}$, leaving a residual
$\bigO(z^{-J-1})$. All relevant derivatives of the finite logarithmic model
are controlled on a bounded neighborhood of $z+d_0$, and its first derivative
is $\log z+\bigO(1/z)$. The actual index is in such a neighborhood: the leading
term first gives $n/z\to1$, and then the mean value theorem applied to
$n\log n-n-z\log z+z=\bigO(\log n)$ gives $n-z=\bigO(1)$.
A further application of the mean value theorem divides the residual by a
quantity asymptotic to $w$, proving \eqref{spf:eq:inverse-all}.
\end{proof}

\subsection{What is being inverted}

A discrete sequence has no distinguished analytic interpolation. The preceding
theorems make no such choice: their statements are evaluated at the actual
values $y=A(n)$. Alternatively, define an explicit smooth approximation
\begin{equation}
 \mathcal A_M(x)=e^{-\theta}\Gamma(x+1)\sum_{j=0}^{M}c_jx^{-j}.
\end{equation}
It is positive and strictly increasing for all sufficiently large $x$, because
the finite correction tends to one and the logarithmic derivative is
$\log x+\bigO(1/x)$. Its inverse satisfies
\begin{equation}
 \mathcal A_M^{-1}(A(n))-n
 =\bigO\left(\frac{n^{-M-1}}{\log n}\right).
\end{equation}
This is a well-defined model inverse, with its discrepancy from the sequence
controlled.

For an integer threshold $N(y)=\min\{n:A(n)\ge y\}$, an asymptotic estimate
alone does not justify unconditional nearest-integer rounding near a jump.
Certified bounds for $A(n)$ can resolve the remaining adjacent candidates.
The Bonferroni method in the next section provides one such tool.
[Added 1 October 2026, batch 73O2: the general form of this remark, exact
rounding $N(y)=\lceil\nu(y)\rceil$ for an increasing interpolation $\nu$ and the
separation condition under which an approximate inverse decides it, is the
staircase theorem \texttt{p0:thm:staircase} of~\cite{tai}. Its separation
step is formalized in Lean as \texttt{Fabius.staircase\_separation} in
\path{Analysis/FabiusFunction/Lean/FabiusFunction/StaircaseInversion.lean}.
That formalization is context only; it does not verify any statement of this
report.]

\section{Reproducible computations and exact certificates}
\label{spf:sec:computation}

\subsection{Coefficient generation}

The core implementation, \texttt{code/path\_forests.py}, uses only the Python
standard library. It forms the numerator factors and reciprocal denominator
in \eqref{spf:eq:Q} as truncated integer polynomials; evaluates
$P_{\boldsymbol\beta,j}(m)$ at the proved sufficient set of integers; takes
forward differences; and assembles \eqref{spf:eq:Bformula} using exact rational
arithmetic. The work depends on the requested asymptotic order, not on a large
sequence index $n$.

At order $J$, profiles are indexed by partitions of $d\le J$, and each profile
uses at most $2(J-d)+1$ dimer evaluations. No integer-sequence terms are used
as fitting data. The archive contains all $B_J$ and $c_J$ through order sixteen
in \texttt{data/coefficients\_order16.json}.

\subsection{Independent finite checks}

The included validation program completed the following checks.
\begin{center}
\small
\begin{tabular}{p{0.65\linewidth}r}
\toprule
Check & Exact equalities/cases\\
\midrule
Full tiling enumeration against both OEIS lists, $n=0,\ldots,21$ & 44\\
Brute-force full distributions for $n=1,\ldots,8$, six offset pairs,
 and both statistics & 96\\
Stable factorial moments against those brute-force distributions & 246\\
Stable tiling polynomial on specified unequal path-length forests & 86\\
Additional finite-difference tests beyond the interpolation degree & 375\\
Factorial denominator checks on polynomial coefficients through order 16 & 396\\
\bottomrule
\end{tabular}
\end{center}
The six offset pairs were $(1,1),(1,2),(2,2),(2,3),(3,2),(3,3)$. Full tiling
enumeration is independent of the stable-polynomial implementation, and brute
permutation enumeration is independent of both. These checks detect finite
implementation and bookkeeping errors; none replaces the proofs for arbitrary
$n$ or arbitrary expansion order.

\subsection{Bonferroni probability intervals}

When \eqref{spf:eq:stable-range} holds through $2q+1$, exact factorial moments give
\begin{equation}
 S_{2q+1}\le\Pp(X=0)\le S_{2q},\qquad
 S_j=\sum_{k=0}^{j}(-1)^k\mu_k.
 \label{spf:eq:bonferroni}
\end{equation}
For completeness, this follows pointwise from
$\sum_{k=0}^{j}(-1)^k\binom{x}{k}=(-1)^j\binom{x-1}{j}$ when $x\ge1$,
while both partial sums equal one for $x=0$. Thus the lower and upper bounds
are rigorous rational numbers, even though the full count has not been
computed.

The certificate program uses $S_{31}$ and $S_{30}$ at $n=80,160,320$.
It brackets $e^{-1}$ between the odd and even Taylor sums of degrees 101 and
102, and squares that interval for $e^{-2}$. All arithmetic in the resulting
error intervals is rational. The quantity displayed below is
\begin{equation}
 E_M(n)=n^{M+1}\left(e^{\theta}\frac{A^{(\theta)}_{2,2}(n)}{n!}
                       -\sum_{j=0}^{M}c_jn^{-j}\right).
 \label{spf:eq:scaled-error}
\end{equation}
By the theorem its limit is $c_{M+1}$.

\begin{table}[htbp]
\centering\small
\renewcommand{\arraystretch}{1.17}
\begin{tabular}{rrll}
\toprule
$\theta$ & $n$ & Certified interval for $E_4(n)$ & Certified interval for $E_8(n)$\\
\midrule
1&80  &$[3.3409847,\ 3.3409848]$ &$[-1584.0080,\ -1584.0079]$\\
1&160 &$[3.1664733,\ 3.1664734]$ &$[-1494.2195,\ -1494.2194]$\\
1&320 &$[3.0822399,\ 3.0822400]$ &$[-1451.0335,\ -1451.0334]$\\
2&80  &$[360.08331,\ 360.08332]$ &$[359523.99,\ 359524.00]$\\
2&160 &$[348.86238,\ 348.86239]$ &$[343097.56,\ 343097.57]$\\
2&320 &$[343.54178,\ 343.54179]$ &$[335473.66,\ 335473.68]$\\
\bottomrule
\end{tabular}
\caption{Outward-rounded exact intervals. The limits are $3,-1409$ for
$\theta=1$, and $1692/5,186091480/567$ for $\theta=2$. The intervals certify
these finite residuals; their convergence is a consequence of the proof, not
an inference from three samples.}
\label{spf:tab:certificates}
\end{table}

The archive includes the unrounded rational endpoints and additional
order-ten residual certificates. With the same moment cutoff, those
order-ten intervals become wider at larger $n$, especially for $\theta=2$:
the factor $n^{11}$ magnifies the fixed Bonferroni width. This is an honest
precision limit of that finite computation; increasing the cutoff addresses
it without using a guessed recurrence.

\subsection{Inverse diagnostics and reproducibility}

Using exact enumerated sequence values at $n=20$, the Lambert backbone $z$
has index errors about $0.52023$ for A189281 and $0.20749$ for A110128.
Including both displayed correction levels in \eqref{spf:eq:inverse-first}
reduces these to approximately $-0.00066084$ and $0.00089419$, respectively.
These are numerical diagnostics, not certified rounding guarantees.

The programs were run with Python 3.13.5. The core coefficient, enumeration,
and rational-certificate programs need no third-party package. The optional
rational-identity exploration uses SymPy 1.14.0; the inverse diagnostic uses
mpmath 1.3.0. The commands and expected outputs are recorded in the archive's
README. No Lean code was executed or supplied as an asserted formal proof.
[Added 1 October 2026, batch 73O2: the delivered README is not shipped. The
report's \texttt{README.md} replaces it and gives rerun commands for the
shipped layout; see \cref{spf:app:provenance}.]

\section{A sharper conjecture suggested by the exact computations}
\label{spf:sec:collapse}

There is a striking simplification in the oriented $(2,2)$ case that is not
needed by any preceding theorem. Define the universal rational moments
\begin{equation}
 \mu_k^*(n)=\sum_{\mathbf b:\,\sum(\ell-1)b_\ell=k}
       \frac{F_2(n;\mathbf b)^2}{\bfac\,\fall nv},
\end{equation}
as rational functions in an indeterminate $n$. They agree with the actual
moments whenever the stability conditions hold. Remove the leading Poisson
factor by setting
\begin{equation}
 R_h(n)=\sum_{k=0}^{h}\frac{(-1)^{h-k}}{(h-k)!}\mu_k^*(n).
 \label{spf:eq:R}
\end{equation}
Exact finite algebra gives
\begin{align}
 R_0&=1,\\
 R_1&=-\frac{3n-4}{\fall n2},\\
 R_2&=\frac{3n^2-17n+26}{\fall n4},\\
 R_3&=-\frac{(n-6)(n^2-7n+16)}{\fall n6}.
\end{align}

\begin{conjecture}[Rational collapse]
\label{spf:conj:collapse}
For every $h\ge4$,
\begin{equation}
 R_h(n)=\frac{24\,\fall{(n-h+1)}{h-4}}{\fall n{2h}}
       =\frac{24}{\fall n{h-1}\,\fall{(n-2h+5)}5}.
 \label{spf:eq:collapse}
\end{equation}
\end{conjecture}

This has been checked as an exact identity of rational functions for each
$h=4,\ldots,16$. The check clears the common denominator $\fall n{2h}$ and
compares polynomials; it is stronger than testing finitely many numerical
values of $n$, but it is still only a finite check in $h$. No induction or
creative-telescoping certificate for all $h$ has been obtained here.

The proposed formula would imply the recurrence
\begin{equation}
 \frac{R_{h+1}(n)}{R_h(n)}
 =\frac{(n-2h+5)(n-2h+4)}{(n-h+1)(n-2h)(n-2h-1)}
 \qquad(h\ge4).
 \label{spf:eq:R-ratio}
\end{equation}
Proving \eqref{spf:eq:R-ratio} directly from \eqref{spf:eq:R} is therefore a sharply
specified next target.

\begin{proposition}[What the rational collapse would imply]
\label{spf:prop:collapse-implication}
If Conjecture~\ref{spf:conj:collapse} holds, then
$B_J^{(1)}(u;2,2)\in\Z[u]$ for every $J$, and therefore all
$c_J^{(1)}(2,2)$ are integers. For $J\ge4$ it would also give
\begin{equation}
 \deg B_J^{(1)}(u;2,2)\le\max(3,J-4).
\end{equation}
\end{proposition}
\begin{proof}
The first four rational functions $R_h(1/x)$ have integer Taylor
coefficients. Under the conjecture, each later one is $24x^{h+4}$ times a
product of factors $(1-ax)^{-1}$ with nonnegative integers $a$, after
cancellation. It too has integer coefficients and begins at degree $h+4$.
At any fixed power $x^J$, only finitely many of these $R_h$ can contribute.

By \eqref{spf:eq:R}, the coefficient of $u^h$ in the Poisson-normalized stable
moment series is $R_h(n)$. Comparing with the already proved coefficient
formula gives
$B_J^{(1)}(u;2,2)=\sum_h[x^J]R_h(1/x)u^h$.
The integrality and degree claims follow.
\end{proof}

This conditional proposition must not be confused with the unconditional
factorial denominator theorem. The conjecture packages the observed
integrality into a more structured rational-function identity, making it
more approachable than a bare statement about an infinite list of integers.

\section{Further research questions and proposed next steps}
\label{spf:sec:questions}

\paragraph{1. Prove or refute the rational collapse.}
The most concrete next step is an identity in two symbolic parameters rather
than a fit to sequence values. Clear denominators in \eqref{spf:eq:R}, organize
profiles by their component count, and seek a telescoping proof of
\eqref{spf:eq:R-ratio}. A combinatorial cancellation proof would be particularly
valuable: it should explain both the factor $24$ and the loss of four powers
of $n$ beyond the moment index. A failure at some later $h$ would also be
informative, because the first discrepancy would identify a missing component
interaction. The all-orders integrality issue for the oriented family remains
separate until such a proof exists.
[Added 1 October 2026, batch 74: for the oriented one-path subfamily,
$r=1$ or $s=1$, all-orders integrality is proved in Part~II
(\cref{fga:cor:integrality}); the case $(2,2)$ of A189281 remains open.]

\paragraph{2. Connect the asymptotic theorem to the specific guessed operators.}
A proved asymptotic expansion does not establish the A189281 or A110128
recurrence. Conversely, a rigorous operator certificate could supply a second,
independent route to large-order coefficient behavior. One useful project is
to translate the stable moment framework into a finite-dimensional holonomic
representation with boundary terms retained, then compare its annihilator
with the listed operators. All boundary terms must be proved, not discarded
because their numerical influence is small.

\paragraph{3. Identify the first geometry-dependent exponential sector.}
Theorem~\ref{spf:thm:universality} shows that long path-length differences are
invisible to every algebraic order. It does not give their leading amplitude,
sign, or possible dependence on $n$ modulo the offsets. The first profiles
that can see a short boundary have edge degree comparable to the shortest
path. A saddle-point analysis of that moving profile boundary could reveal
the first exponentially small sector, and would turn the present algebraic
expansion into part of a genuine transseries description.

\paragraph{4. Determine large-order growth and justified optimal truncation.}
The coefficients through order sixteen already show that sign and magnitude
can change markedly. The present proof allows constants to depend on $M$.
It would be useful to prove explicit growth bounds in $M$, decide whether a
Gevrey-type estimate is sharp, and derive a valid truncation rule for a given
$n$. The rational-collapse conjecture, if proved, would make the oriented
case a tractable testing ground, but the absolute case needs its own analysis.

\paragraph{5. Allow offsets or path counts to grow with $n$.}
The current proof uses fixed $r,s$, both in the shortest-path condition and in
the analytic bound on $Q_{\mathbf b}$. When $r$ or $s$ grows, the shortest path
may meet the moment cutoff, and the number of path boundaries changes the
first corrections. Determine the largest growth range in which a useful
uniform expansion persists, and the transition law when path lengths become
comparable to $\log n/\log\log n$. This question requires new uniform estimates;
it is not an immediate corollary of the fixed-offset theorem.
[Added 1 October 2026, batch 74: for fixed $r,s$, Part~II shows that the
expansion survives whenever the shortest path is long compared with
$\log n/\log\log n$ (\cref{fga:cor:shortest}); the transition at that scale
remains open (\cref{fga:sec:questions}, question~2).]

\paragraph{6. Extend from paths to several forbidden offsets and to cycles.}
For several allowed differences, selected components can branch or contain
cycles. Their embeddings are no longer determined by one starting value and a
single orientation. A useful extension would classify component types by an
appropriate defect and prove an analogue of the uniform moment bound and
finite coefficient formula. Cyclic boundaries are another natural test:
short cycles can contribute at a finite algebraic order, whereas long cycles
may only change a beyond-all-orders sector.

\paragraph{7. Formalize the proof and build a certified coefficient library.}
The finite profile identity, the stable tiling polynomial, and the integer
coefficient arithmetic are natural first formalization targets. The analytic
layer then needs the uniform factorial-moment tail, Cauchy bounds, and finite
asymptotic truncation. A robust implementation should expose the exact
stability hypotheses instead of silently applying $F_r$ outside its range.
The rational Bonferroni certificates are suitable artifacts for an independent
checker. This is a proposed formalization route, not a report of completed
Lean verification.

\paragraph{8. Combine asymptotic inversion with exact threshold certification.}
For very large $y$, the Lambert backbone and a few correction terms should
reduce the inverse problem to a small number of candidate indices. Use exact
Bonferroni intervals to compare each candidate count with $y$, adaptively
increasing the moment cutoff only when necessary. The interesting algorithmic
question is a proved cost bound that accounts for near-threshold inputs,
where asymptotic rounding alone cannot decide the answer.

\section{Conclusion}

The main result is a direct all-orders theorem for a natural pair of
permutation sequences whose higher displayed OEIS expansions were tied to
guessed recurrences. Its proof does not solve those recurrence problems by
assumption: it bypasses them. The stable path-tiling identity gives exact low
moments, the factorial tail controls the remaining moments, and finite
differences sum the unlimited dimer contributions. Together they turn a formal
asymptotic calculation into a uniform theorem.

The same structure yields more than two lists of numbers: a full-law
Poisson correction expansion, a quantitative explanation of boundary-length
universality, an arithmetic denominator theorem, and a rigorous inverse
construction. The stronger rational-collapse pattern is separated from these
proved results and posed as a specific next problem. That distinction is
central to the usefulness of the report: the established theorems can be
used now, while the remaining conjectures have explicit, testable targets.

\part{Fixed-gap asymptotics: a second derivation and its own results}
\label{fga:part}

\section{Provenance, scope and notation of Part II}
\label{fga:sec:front}

[Added 1 October 2026, batch 74.] Part~II was added in the write of
batch~74. It is built from manuscript~03 of that batch and prints only what
that manuscript adds to Part~I. Where the two manuscripts prove the same
statement, the statement is printed once, in Part~I, and manuscript~03's
route is recorded in \cref{fga:sec:routes}. Everything in Part~II is
written in Part~I's notation (\cref{fga:tab:dictionary}).

\subsection{Provenance}
\label{fga:sec:provenance}

Manuscript~03 of batch~74 is the archive
\nolinkurl{ProveIt_Fixed_Gap_Asymptotics.zip}, which arrived in commit
\texttt{c664fc2f0} and was placed in this report by commit
\texttt{b669cff87}. Its article is titled \emph{All-Orders Asymptotics and
Boundary Universality for Fixed-Gap Permutations} (25~pages, dated
October~1, 2026); its title page reads ``Prepared for Vladimir
Reshetnikov'' and its PDF metadata names ``ChatGPT, research report
prepared for Vladimir Reshetnikov'' as author. Its pinned repository commit
is \texttt{291fbe6bb4477e0cb2131f1e041b12a9abe363fd} (2026-10-01, 18:10;
\cite{fgaProveIt}).

It was written without sight of manuscript~60, the source of Part~I, and
proves Part~I's main theorems in independent text: its structure,
notation, constants and several proof devices differ throughout, and the
two texts share about 1.1\,\% of their word 8-grams, almost all of them
bibliography entries and the digits of the coefficient list. Every number
the two manuscripts share agrees exactly: the coefficients of A189281 and
A110128 through order twelve, the general second correction in $r,s$ for
both orientations, the polynomials $B_1,B_2,B_3$ at $(r,s)=(2,2)$, the
logarithmic coefficients and the first inverse coefficients.

[Added 1 October 2026, batch 74: manuscript~03 states that ``a repository
code search for A189281 returned no match at the time of the audit'', and
its delivered \texttt{02-fixed-gap-sources.md} repeats it. The pinned
commit \texttt{291fbe6bb} already held the archive of manuscript~60,
\texttt{docs/incoming/oeis\_path\_forests.zip} (arrival commit
\texttt{aa43cc555}), which code search cannot read; that archive was placed
an hour later as Part~I of this report (commit \texttt{6e193dd4f}). Part~II
is an independent derivation of the same theorems plus the results listed
below.]

\paragraph{Scope and non-claims of manuscript~03, kept.} Its contribution
is ``the explicit stable-profile and finite-defect framework, the uniform
error analysis, and the quantitative comparison of whole laws''. The
leading Poisson behaviour, Tauraso's first correction and the
matching-of-tilings inclusion--exclusion of Spahn and Zeilberger are
credited as prior work. The proofs are ordinary mathematics, not a
Lean-certified development, and no publication priority is claimed. ``All
orders'' means a remainder theorem at each fixed order, not a convergent or
resurgent completion; no Borel summability, optimal truncation, Stokes
constant or complete family of exponentially small sectors is asserted. The
conjectured OEIS recurrences of A189281 and A110128 are not proved, and
neither is all-orders integrality when both component counts exceed one.
Its numerical tables are not interval certificates, and its values of
A189281 at $n=80,120$ are b-file inputs, not regenerated. It made no OEIS
submission and modified no repository file.

\subsection{What Part II adds, and how it compares with Part I}
\label{fga:sec:claims}

\begin{table}[htbp]
\centering\small
\renewcommand{\arraystretch}{1.15}
\begin{tabular}{@{}>{\raggedright\arraybackslash}p{.2\textwidth}>{\raggedright\arraybackslash}p{.36\textwidth}>{\raggedright\arraybackslash}p{.36\textwidth}@{}}
\toprule
Topic & Part I (manuscript 60) & Part II (manuscript 03)\\
\midrule
All-orders remainder &
fixed offsets $r,s$; uniform on every disk $|u|\le U$
(\cref{spf:thm:main}) &
only on $|1+u|\le1$, i.e.\ $|u|\le2$, but uniform over all forest shapes
whose shortest path has at least $K_n\asymp\log n/\log\log n$ vertices
(\cref{fga:thm:allorders,fga:cor:shortest})\\
Coefficients &
$B_J$ through order 16, $(2J)!B_J\in\Z[u]$ (Proposition~\ref{spf:prop:arithmetic}) &
same values through order 12 (not new); $B_J$ polynomial in $r,s$ of total
degree at most $J$ (\cref{fga:prop:rs-degree}); closed forms of
$c_3^{(\theta)}(r,s)$ (\cref{fga:prop:c3})\\
Integrality &
for $(2,2)$ oriented only conditional on Conjecture~\ref{spf:conj:collapse} &
\emph{proved} for the oriented case with $r=1$ or $s=1$
(\cref{fga:cor:integrality}); $(2,2)$ untouched\\
Universality &
total variation bound (\cref{spf:thm:universality}) &
bounds for the whole generating function and the avoidance probability,
a second total-variation bound, and an $n=12$ example
(\cref{fga:thm:universal,fga:rem:n12})\\
Whole law &
$\ell^1$ approximation of the law (Corollary~\ref{spf:cor:distribution}) &
explicit first-order local laws at $(2,2)$ \eqref{fga:eq:localdirect}--\eqref{fga:eq:localabsolute}\\
Inverse &
all orders in the coordinates $z,w$ (\cref{spf:thm:inverse-all}) &
explicit terms through $N^{-3}$ in a coordinate absorbing $\kappa_\theta$;
$a_2,a_3$ for A110128 (\cref{fga:prop:inverse})\\
Certificates &
certified Bonferroni intervals (\cref{spf:tab:certificates}) &
uncertified decimals only (\cref{fga:sec:numerics})\\
\bottomrule
\end{tabular}
\caption{Part II against Part I. The remainder theorem of Part~II is weaker
in $u$ and stronger in geometry; see also \cref{fga:rem:partI-shapes}.}
\label{fga:tab:claims}
\end{table}

\subsection{Notation}
\label{fga:sec:notation}

Manuscript~03 and Part~I use several letters for different objects, and
manuscript~03 overloads $D$, $w$, $T$, $A$ and $B$ internally.
\Cref{fga:tab:dictionary} fixes the convention: Part~II uses Part~I's names
for every shared object and new names ($K_n$, $D_{n,k}$, $\Lambda$, $N$,
$\Delta$, $\omega_j$, $\varepsilon_M$, $\alpha$) for the rest. No
normalization is changed.

\begin{table}[htbp]
\centering\small
\renewcommand{\arraystretch}{1.12}
\begin{tabular}{@{}>{\raggedright\arraybackslash}p{.33\textwidth}>{\raggedright\arraybackslash}p{.27\textwidth}>{\raggedright\arraybackslash}p{.32\textwidth}@{}}
\toprule
Notion & Part I and Part II & Manuscript 03\\
\midrule
orientation weight ($1$ or $2$) & $\theta$ & $\eta$ (Part~I's $\eta$ is a
length proportion)\\
length proportion & $\alpha$ (Part~II); $\eta$ in \cref{spf:thm:universality} & $\alpha$\\
generating function & $G(1+u)=\E(1+u)^X$ & $G(w)$, $w=1+z$, $z=u$\\
$\E\binom Xk$ & $\mu_k$ & $S_k$ (Part~I's $S_j$ are Bonferroni sums)\\
profile, multiplicity factorial & $\mathbf b=(b_\ell)$, $\bfac$ & $\mathbf m=(m_j)$, $\mathbf m!$\\
tiling count & $C_{\mathcal P}(\mathbf b)$ & $T_{\mathcal F}(\mathbf m)$\\
stable tiling polynomial & $F_r(n;\mathbf b)$ & $\bfac^{\,-1}F_r$ is its $T_r(n;\mathbf m)$\\
formal root & $1+W$, $\mathcal D^{-1}$ & $\lambda$, $A$\\
dimers; long profile & $m$; $\boldsymbol\beta$, $C$, $K$ & $h$; $\mathbf m$, $c_0$, $k_0$\\
normalized profile kernel & $Q_{\mathbf b}(x)$ & $H_h(t)$ (\emph{not} Part~I's $H_h$)\\
Charlier polynomials & $H_h(q;\theta)$ & $Q_m(j;\eta)$ (\emph{not} Part~I's $Q_{\mathbf b}$)\\
correction polynomials & $B_J^{(\theta)}(u;r,s)$ & $C_J(r,s;\eta,z)$\\
avoidance coefficients & $c_J^{(\theta)}(r,s)$ & $c_J^{(\eta)}(r,s)$ (same)\\
moment cutoff & $T_n$ (Part~I), $K_n$ (Part~II) & $K$ (Part~I's $K$ is an edge degree)\\
$(1-2k/n)^{-2k}$ & $D_{n,k}$ & $D_{n,k}$ (its $D$ is also $\log N$)\\
logarithmic coefficients & $a_j$ & $d_j$ (Part~I's $d_j$ are inverse coefficients)\\
$\tfrac12\log(2\pi)-\theta$ & $\kappa_\theta$ & $\gamma$\\
inverse coordinates & $\Lambda$, $N$, $\Delta=\log N$ (Part~II) & $L$, $N$, $D$\\
inverse coefficients & $\omega_j(\Delta)$ & $u_j(D)$\\
unscaled truncation error & $\varepsilon_M(n)$ & $E_M(n)$ (Part~I's $E_M$ is scaled by $n^{M+1}$)\\
\bottomrule
\end{tabular}
\caption{Notation of manuscript~03 against Part~I. Part~II uses the middle
column throughout.}
\label{fga:tab:dictionary}
\end{table}

\section{Manuscript 03's second routes to Part I's theorems}
\label{fga:sec:routes}

\Cref{fga:tab:routes} records, for each result of Part~I that
manuscript~03 proves again, the device its proof uses. None of these
statements is reprinted. Where a device is also needed for a new result
of Part~II, it is stated there.

\begin{table}[htbp]
\centering\small
\renewcommand{\arraystretch}{1.12}
\begin{tabular}{@{}>{\raggedright\arraybackslash}p{.27\textwidth}>{\raggedright\arraybackslash}p{.66\textwidth}@{}}
\toprule
Part I & Manuscript 03's route\\
\midrule
Proposition~\ref{spf:prop:exact} &
the same matching-of-tilings argument, stated for binomial moments $\mu_k$,
with the same orientation lemma (weight $\theta^c$, not $\theta^k$); adds
the symmetry $A_{r,s}=A_{s,r}$ via $\pi\mapsto\pi^{-1}$\\
\eqref{spf:eq:lagrange-id} &
a formal residue identity for the root $\lambda=1+\sum_\ell
z_\ell\lambda^{1-\ell}$, with $A=(1+\sum(\ell-1)z_\ell\lambda^{-\ell})^{-1}$;
here $\lambda=1+W$ and $A=\mathcal D^{-1}$\\
\cref{spf:thm:stable} &
the identical polynomial, written $T_r(n;\mathbf m)=F_r(n;\mathbf b)/\bfac$\\
Lemma~\ref{spf:lem:moment-bound} &
$\mu_k\le\theta^kD_{n,k}/k!$ for $1\le k<n/2$ (\cref{fga:lem:momentbound}):
sharper for $k\ll n$, Part~I's holds for every $k$\\
\eqref{spf:eq:moment-tail} &
a Taylor remainder on the unit disk $|1+u|\le1$ (\cref{fga:lem:taylor}),
not an exponential-series tail\\
\cref{spf:thm:universality} &
total variation through Parseval and Cauchy--Schwarz
(\cref{fga:thm:universal}); the $\exp(-\delta n\log n+O(n))$ form is
Part~I's last assertion with $\eta'$ renamed $\delta$\\
Definition~\ref{spf:def:B} &
the same finite formula; the dimer sum is evaluated by Touchard polynomials,
$e^{-x}\sum_mx^mm^q/m!=T_q(x)$ with $T_{q+1}=x(T_q+T_q')$, instead of
finite differences, and the denominator by exponentials of power sums\\
Proposition~\ref{spf:prop:arithmetic} &
$B_0=1$, $\deg_uB_J\le2J$, $B_J(0)=0$ and symmetry; adds the degree in
$r,s$ (\cref{fga:prop:rs-degree}); Part~I's denominator theorem is not in
manuscript~03\\
\eqref{spf:eq:B1}, \eqref{spf:eq:B2-oriented}, \eqref{spf:eq:B2-absolute} &
$B_1$ by the same two-profile hand calculation; general $B_2$ and $B_3$ in
symbolic $r,s$ in \nolinkurl{data/02-fixed-gap-general_pgf_first_three.txt};
$B_2$ agrees with Part~I identically in $r,s,u$\\
\cref{spf:thm:main} &
\cref{fga:thm:allorders}: a defect-sensitive moment bound and a kernel
bound on $|x|\le1/(D(1+k)^2)$ in place of Lemma~\ref{spf:lem:Qbound}\\
\eqref{spf:eq:A189281}, \eqref{spf:eq:A110128}, \cref{spf:tab:extra} &
the coefficients of both sequences through order~12, equal to Part~I's
(see below)\\
residue classes (\cref{spf:sec:target}) &
the same explanation by \cref{spf:thm:stable}\\
\eqref{spf:eq:exact-mean} &
the same count; adds
$\E X=\theta+\theta(1-r-s)/n+\theta(r-1)(s-1)\sum_{q\ge2}n^{-q}$\\
\eqref{spf:eq:logA}, \cref{spf:thm:inverse-first,spf:thm:inverse-all} &
the same logarithmic series; the inverse in a coordinate absorbing
$\kappa_\theta$ (\cref{fga:prop:inverse})\\
\cref{spf:sec:inverse}, last subsection &
the same smooth carrier $\mathcal A_M$ and threshold bracket\\
\bottomrule
\end{tabular}
\caption{Results of Part~I proved again by manuscript~03, and the device of
each second route.}
\label{fga:tab:routes}
\end{table}

\paragraph{Order twelve is not new.} Manuscript~03 presents its
coefficients through order twelve as an extension of the OEIS lists
(``extend them through order twelve'') and proposes appending
\[
 -\frac{59103}{n^{11}}-\frac{9448}{n^{12}}\ \text{to A189281},\qquad
 \frac{2620164054016}{155925\,n^{11}}+\frac{61699558213516}{467775\,n^{12}}
 \ \text{to A110128}.
\]
[Added 1 October 2026, batch 74: Part~I's \cref{spf:tab:extra} already
gives orders 11 to~16 of both sequences. Manuscript~03's values at orders
0 to~12, computed independently by its own engine, agree with Part~I's
exactly; they confirm Part~I's table and are not new.]

\paragraph{The inverse is not new as a method.} As for Part~I
(Remark~\ref{spf:rem:inverse-apparatus}), manuscript~03's inverse is an instance
of the transseries volume's apparatus~\cite{tai}: its coordinates are those
of \texttt{p6:thm:gamma} with $\log Y$ shifted by $\theta$, its inverse
expansion is a perturbed inversion around that gamma carrier
(\texttt{p0:thm:perturbed-inversion}), and its threshold bracket is
\texttt{p0:thm:staircase}. No novelty is claimed for the method.

\section{Uniformity over forest shapes}
\label{fga:sec:shapes}

Throughout this section $\mathcal P$ and $\mathcal V$ are arbitrary
directed path forests on $n$ vertices with $r$ and $s$ paths, as in
\cref{spf:sec:tilings}, and
\begin{equation}
 D_{n,k}=\Bigl(1-\frac{2k}n\Bigr)^{-2k}\qquad(1\le k<n/2).
 \label{fga:eq:Dnk}
\end{equation}

\begin{lemma}[A sharper binomial-moment bound; manuscript~03]
\label[lemma]{fga:lem:momentbound}
For either statistic and arbitrary path forests,
$\mu_k\le\theta^kD_{n,k}/k!$ for $1\le k<n/2$.
\end{lemma}
\begin{proof}[Proof (manuscript~03)]
As in Lemma~\ref{spf:lem:moment-bound}, a selected set of $k$ position edges
with $c$ components has probability at most $\theta^cn^c/\fall n{k+c}$. Since
$c\le k$ and $k+c\le2k$, this is at most $\theta^kn^{-k}D_{n,k}$; there are
at most $n^k/k!$ such sets.
\end{proof}

\begin{lemma}[Taylor remainder on the unit disk; manuscript~03]
\label[lemma]{fga:lem:taylor}
For $|1+u|\le1$ and every integer $K\ge1$,
\begin{equation}
 \Bigl|G(1+u)-\sum_{k=0}^{K-1}\mu_ku^k\Bigr|\le|u|^K\mu_K.
 \label{fga:eq:taylor}
\end{equation}
\end{lemma}
\begin{proof}[Proof (manuscript~03)]
For an integer $x\ge K$, Taylor's integral remainder for $w^x$ about $w=1$
along the segment from $1$ to $w=1+u$, which stays in the closed unit disk,
has modulus at most $\binom xK|u|^K$; for $x<K$ it vanishes. Take
expectations.
\end{proof}

\begin{theorem}[Quantitative boundary universality; manuscript~03]
\label{fga:thm:universal}
Let $(\mathcal P,\mathcal V)$ and $(\widetilde{\mathcal P},\widetilde{\mathcal V})$
be two pairs of forests on $n$ vertices with $r$ position and $s$ value
paths, every path of all four having at least $K$ vertices, $1\le K<n/2$,
and let $X,\widetilde X$ be their statistics for the same $\theta$. Then
$\mu_k=\widetilde\mu_k$ for $0\le k\le K$, and
\begin{align}
 \sup_{|w|\le1}\bigl|G(w)-\widetilde G(w)\bigr|
 &\le\frac{2(2\theta)^K}{K!}D_{n,K},\label{fga:eq:universalpgf}\\
 \bigl|\Pp(X=0)-\Pp(\widetilde X=0)\bigr|
 &\le\frac{2\theta^K}{K!}D_{n,K},\label{fga:eq:universalavoid}\\
 \TV\bigl(\mathcal L(X),\mathcal L(\widetilde X)\bigr)
 &\le\sqrt{n+1}\,\frac{(2\theta)^K}{K!}D_{n,K}.\label{fga:eq:universaltv}
\end{align}
If every path has at least $\alpha n$ vertices, all three are at most
$\exp\{-\delta n\log n+O_{\alpha,\delta,\theta}(n)\}$ for every fixed
$0<\delta<\min(\alpha,1/2)$.
\end{theorem}
\begin{proof}[Proof (manuscript~03)]
Moments of order $k\le K$ are stable in all four forests and agree by
Corollary~\ref{spf:cor:stable-moments}. Subtract the common Taylor polynomial of
degree $K-1$ and apply \cref{fga:lem:taylor,fga:lem:momentbound} to both
remainders, with $|u|=|w-1|\le2$ on the disk and $|u|=1$ at $w=0$. For the
total variation, write $G-\widetilde G=\sum_{j\le n}q_jw^j$; Parseval on
the unit circle gives $\sum|q_j|^2\le\sup_{|w|=1}|G-\widetilde G|^2$, and
Cauchy--Schwarz gives
$\frac12\sum|q_j|\le\frac12\sqrt{n+1}\sup|G-\widetilde G|$. For the last
assertion take $K=\lfloor\delta n\rfloor$: $\log D_{n,K}=O(n)$ and
$\log K!=\delta n\log n+O(n)$.
\end{proof}

The bounds \eqref{fga:eq:universalpgf}--\eqref{fga:eq:universalavoid} on
the whole generating function and on the avoidance probability are new
here; \eqref{fga:eq:universaltv} and the $\exp(-\delta n\log n)$ form are a
second route to \cref{spf:thm:universality}. Boundary placement can change
the exact count, but, when all components are long, its normalized
influence is smaller than every power of $1/n$; the theorem compares
shape-dependent corrections and does not identify an exponentially small
expansion for a single sequence.

\begin{remark}[Agreement of moments is not equality of laws; manuscript~03]
\label[remark]{fga:rem:n12}
At $n=12$ compare the position and value path lengths $(6,6)/(6,6)$ with
$(5,7)/(4,8)$. Their moments agree through $k=4$ and first differ at $k=5$.
The oriented avoidance counts are $222822578$ and $222822556$, the absolute
ones $91310696$ and $91284560$, by manuscript~03's independent exact tiling
routine. ``All algebraic orders agree'' must not be read as ``the finite
distributions are identical''.
\end{remark}

\begin{lemma}[A defect-sensitive moment bound; manuscript~03]
\label[lemma]{fga:lem:defect}
Let $\mu_{k,c}$ be the part of $\mu_k$ from selected sets with exactly $c$
components, so $d=k-c$. For $1\le c\le k<n/2$,
\begin{equation}
 \mu_{k,c}\le\frac{\theta^c}{c!}\binom{k-1}{c-1}n^{-d}D_{n,k}
 \le\frac{\theta^k}{k!}\Bigl(\frac{k^2}n\Bigr)^dD_{n,k}.
 \label{fga:eq:defect}
\end{equation}
\end{lemma}
\begin{proof}[Proof (manuscript~03)]
There are at most $\frac{n^c}{c!}\binom{k-1}{c-1}$ selected sets with $c$
components and $k$ edges: choose the $c$ initial vertices, order them, and
give them a composition of $k$ into $c$ positive edge counts (overlaps
are overcounted). Each has probability at most
$\theta^cn^c/\fall n{k+c}\le\theta^cn^{-k}D_{n,k}$, because
$\fall n{k+c}\ge n^{k+c}(1-2k/n)^{2k}$ as in \cref{fga:lem:momentbound}.
Multiplying, and using $n^c\cdot n^{-k}=n^{-d}$, gives the first bound.
For the second,
$\binom{k-1}{c-1}=\binom{k-1}d\le k^d$, $k!/c!\le k^d$ and $\theta\ge1$.
\end{proof}

\begin{theorem}[All orders, uniformly over forest shapes; manuscript~03]
\label{fga:thm:allorders}
Fix $r,s\ge1$, $\theta\in\{1,2\}$ and $M\ge0$, and for large $n$ put
\begin{equation}
 K_n=\Bigl\lfloor\frac{(M+4)\log n}{\log\log n}\Bigr\rfloor.
 \label{fga:eq:Kn}
\end{equation}
For every pair of path forests on $n$ vertices with $r$ and $s$ paths, each
path having at least $K_n$ vertices,
\begin{equation}
 G(1+u)=e^{\theta u}\sum_{J=0}^{M}\frac{B_J^{(\theta)}(u;r,s)}{n^J}
 +\bigO_{M,r,s}(n^{-M-1}),\qquad|1+u|\le1,
 \label{fga:eq:allorders}
\end{equation}
with $B_J$ the polynomials of \eqref{spf:eq:Bformula}. The implied constant
does not depend on the individual path lengths.
\end{theorem}
\begin{proof}[Proof (manuscript~03, condensed)]
\emph{Tail.} By \cref{fga:lem:taylor,fga:lem:momentbound} with $|u|\le2$,
truncating at $k<K_n$ costs at most $(2\theta)^{K_n}D_{n,K_n}/K_n!$. Since
$K_n^2/n\to0$, $D_{n,K_n}=\exp(\bigO(K_n^2/n))$, and
$\log K_n!=(M+4+\smallo(1))\log n$; so the cost is $\smallo(n^{-M-1})$. Every
retained profile has $k<K_n$ edges and is stable in both forests
(Corollary~\ref{spf:cor:stable-moments}).

\emph{High defect.} By \cref{fga:lem:defect}, with $k^2/n<1/2$ for $k<K_n$,
the retained profiles of defect $d\ge M+1$ contribute at most
$\bigO(n^{-M-1})\sum_{k\ge1}(2\theta)^kk^{2M+2}/k!=\bigO_M(n^{-M-1})$,
uniformly in $u$ and in the path lengths.

\emph{Low defect.} For each of the finitely many long profiles
$\boldsymbol\beta$ of defect $d\le M$, expand the normalized kernel
$Q_{(m,\boldsymbol\beta)}(1/n)$ of \eqref{spf:eq:Q} through order $M-d$.
Manuscript~03 proves that the kernel is analytic and bounded by a constant
depending only on $r,s$ on $|x|\le1/(D_0(1+k)^2)$ for a suitable $D_0$,
which is the role of Lemma~\ref{spf:lem:Qbound}; Cauchy's estimate gives a
remainder $\bigO((1+k)^{2(M-d)+2}n^{-(M-d)-1})$, summable against the
weights $(\theta|u|)^m/m!$ of the dimers.

\emph{Completion.} Extending each retained dimer sum from $m+K<K_n$ to all
$m\ge0$ adds tails of order $\exp(-K_n\log K_n+\bigO(K_n))=\smallo(n^{-M-1})$.
The completed sums are the Poisson transforms \eqref{spf:eq:poisson-transform}
of the polynomials \eqref{spf:eq:P}, and collecting the powers of $1/n$
gives \eqref{spf:eq:Bformula}. No step uses the path lengths beyond the
stability of the retained profiles.
\end{proof}

\begin{corollary}[How long is long enough; manuscript~03]
\label[corollary]{fga:cor:shortest}
For fixed $r,s$, the expansion \eqref{fga:eq:allorders} holds to every fixed
order for every family of forests whose shortest path length
$L_{\min}(n)$ satisfies
\begin{equation}
 \frac{L_{\min}(n)\log\log n}{\log n}\longrightarrow\infty.
 \label{fga:eq:weaklength}
\end{equation}
In particular the paths need not have lengths proportional to $n$.
\end{corollary}
\begin{proof}[Proof (manuscript~03)]
For each fixed $M$, \eqref{fga:eq:weaklength} eventually gives
$L_{\min}(n)\ge K_n$.
\end{proof}

\begin{remark}[Part~I's proof gives the same uniformity on every disk]
\label[remark]{fga:rem:partI-shapes}
[Added 1 October 2026, batch 74; this observation is the write's, not
either manuscript's.] The remark after the proof of \cref{spf:thm:main}
claims the theorem for forests whose shortest path is a positive multiple
of $n$. Its proof uses the path lengths only to know that every retained
profile, of edge degree at most
$T_n=\lfloor(2M+4)\log n/\log\log n\rfloor$, is stable: the moment tail
\eqref{spf:eq:moment-tail}, the kernel bound Lemma~\ref{spf:lem:Qbound} and
the defect majorant \eqref{spf:eq:defect-majorant} hold for arbitrary
forests. Read so, it gives \cref{spf:thm:main} uniformly on $|u|\le U$ for
all forests whose shortest path has at least $T_n$ vertices, which also
covers the condition \eqref{fga:eq:weaklength}. Manuscript~03 states and
proves the shape-uniform form, on the unit disk only.
\end{remark}

The theorem gives a remainder at each fixed order; it does not assert that
the series in $1/n$ converges, and its proof provides no small optimized
numerical constant for a given finite $n$.

\section{The correction polynomials as polynomials in the offsets}
\label{fga:sec:offsets}

\begin{proposition}[Degree in the offsets; manuscript~03]
\label[proposition]{fga:prop:rs-degree}
For every $J$ and $\theta$ there is a polynomial in $\Q[r,s,u]$ of total
degree at most $J$ in $(r,s)$ whose values at positive integers $r,s$ are
$B_J^{(\theta)}(u;r,s)$; it is symmetric in $r,s$ and has degree at most
$2J$ in $u$.
\end{proposition}
\begin{proof}[Proof (manuscript~03, in Part~I's notation)]
In \eqref{spf:eq:Q}, the offset $r$ enters $f_{r,\mathbf b}(x)$ only
through the rising factorials $\rise{(r-1)}h$, of degree $h$ in $r$, and
each is accompanied by $x^h$; so $[x^j]f_{r,\mathbf b}$ has degree at most
$j$ in $r$, and likewise for $s$. Hence
$P_{\boldsymbol\beta,j}=[x^j]Q_{(m,\boldsymbol\beta)}$ has total degree at
most $j$ in $(r,s)$, and \eqref{spf:eq:Bformula} uses
$P_{\boldsymbol\beta,J-d}$ with $d\ge0$. The other assertions are
Proposition~\ref{spf:prop:arithmetic}.
\end{proof}

\begin{proposition}[The third avoidance coefficient; manuscript~03]
\label[proposition]{fga:prop:c3}
For all positive integers $r,s$,
\begin{align}
 c_1^{(1)}&=r+s-1, & c_1^{(2)}&=2(r+s-2),\notag\\
 c_2^{(1)}&=\tfrac12(r^2-r+s^2-s), & c_2^{(2)}&=2(r^2+rs-2r+s^2-2s),\notag\\
 c_3^{(1)}&=\tfrac16(r+s-1)(r^2-4rs+4r+s^2+4s-6),\label{fga:eq:c3}\\
 c_3^{(2)}&=\tfrac23\bigl(2r^3+3r^2+9rs-17r+2s^3+3s^2-17s+10\bigr).\notag
\end{align}
\end{proposition}

The first four are Part~I's \eqref{spf:eq:avoidance-expansion},
\eqref{spf:eq:B2-oriented} and \eqref{spf:eq:B2-absolute} at $u=-1$; the
closed forms of $c_3^{(\theta)}(r,s)$ are new. Manuscript~03 obtains them by
evaluating its finite formula symbolically in $r,s$ (the full $B_3$ for both
orientations is in \texttt{data/02-fixed-gap-general\_pgf\_first\_three.txt});
this write checked both against Part~I's generator at
$(r,s)=(1,1),(1,3),(2,3),(3,3),(2,5)$. At $r=s=d$, $c_1^{(2)}=4(d-1)$ is
Tauraso's correction \eqref{spf:eq:tauraso}.

\section{The oriented one-path family: proved integrality}
\label{fga:sec:onepath}

\begin{proposition}[Rational collapse for one path; manuscript~03]
\label[proposition]{fga:prop:onepath}
Let the position forest be a single path ($r=1$), the value forest have $s$
paths, and $\theta=1$. Then
\begin{equation}
 n!\,\mu_k=\binom{n-s}k(n-k)!\qquad(0\le k\le n-s),
 \label{fga:eq:onepathmoment}
\end{equation}
so that
$G(1+u)=\frac1{\fall ns}\sum_{k=0}^{n-s}\frac{u^k}{k!}\fall{(n-k)}s$, and,
for $|u|\le2$ and a suitable constant $C_s$,
\begin{equation}
 G(1+u)=\frac{e^u}{\fall ns}\sum_{t\ge0}\frac{\Delta^tp_n(0)}{t!}u^t
 +\bigO\bigl(\exp\{-n\log n+C_sn\}\bigr),
 \qquad p_n(m)=\fall{(n-m)}s.
 \label{fga:eq:onepathrational}
\end{equation}
For $s=1$ the rational factor is $1-u/n$; for $s=2$ it is
$1-2u/n+u^2/(n(n-1))$.
\end{proposition}
\begin{proof}[Proof (manuscript~03)]
Work on the value side. Any $k$ of its $n-s$ edges contract the value set
to $n-k$ blocks, and these blocks can be laid along the position path in
$(n-k)!$ ways, with no further restriction; this is
\eqref{fga:eq:onepathmoment}. Then use
$\fall{(n-s)}k/\fall nk=\fall{(n-k)}s/\fall ns$. Completing the sum to all
$k\ge0$ adds only zero terms through $k=n$ and a factorial tail bounded by
a degree-$s$ polynomial times $2^k/k!$, divided by $\fall ns$; the completed
sum is $e^u$ times the finite-difference Poisson transform
\eqref{spf:eq:poisson-transform} of $p_n$, a polynomial of degree $s$ in
$m$.
\end{proof}

\begin{corollary}[An integrality subfamily; manuscript~03]
\label[corollary]{fga:cor:integrality}
In the oriented case, if $r=1$ or $s=1$, then
$B_J^{(1)}(u;r,s)\in\Z[u]$ for every $J$, and every avoidance coefficient
$c_J^{(1)}(r,s)$ is an integer.
\end{corollary}
\begin{proof}[Proof (manuscript~03)]
The finite sum in \eqref{fga:eq:onepathrational} is
$\sum_qa_q(n)T_q(u)$, where $p_n(m)=\sum_qa_q(n)m^q$ has integer
coefficients and $T_q(u)=e^{-u}\sum_{m\ge0}u^mm^q/m!$ are the Touchard
polynomials, which have integer coefficients; so it is an integer
polynomial in $n$ and $u$ of degree $s$ in $n$. Dividing by
$\fall ns=n^s\prod_{i<s}(1-i/n)$ and expanding
$\prod_{i<s}(1-i/n)^{-1}$, whose power-series coefficients are integers,
gives an expansion in $1/n$ with coefficients in $\Z[u]$. By uniqueness of
asymptotic coefficients it is $\sum_JB_J^{(1)}(u;1,s)n^{-J}$. The case
$s=1$ follows by the symmetry $B_J(u;r,s)=B_J(u;s,r)$.
\end{proof}

This proves integrality only for the oriented one-path subfamily. It says
nothing about the case $(2,2)$ of A189281, whose all-orders integrality
remains conditional on Conjecture~\ref{spf:conj:collapse}
(Proposition~\ref{spf:prop:collapse-implication}); manuscript~03 likewise warns that
the subfamily does not justify extrapolating convergence, termination or
integrality to two paths on both sides. At $(r,s)=(1,1)$ the factor
$1-u/n$ gives $c_J^{(1)}(1,1)=1,1,0,0,\ldots$, the values found in
Remark~\ref{spf:rem:offsets-one-one}.

\section{Fixed local probabilities}
\label{fga:sec:local}

For fixed $j$ and $M$, coefficient extraction from \eqref{spf:eq:main} on a
fixed circle gives
\begin{equation}
 \Pp(X=j)=e^{-\theta}\frac{\theta^j}{j!}\Bigl\{\sum_{J=0}^M\frac1{n^J}
 \sum_hb_{J,h}H_h(j;\theta)+\bigO_j(n^{-M-1})\Bigr\},
 \label{fga:eq:localall}
\end{equation}
with $B_J=\sum_hb_{J,h}u^h$ and the Charlier polynomials
\eqref{spf:eq:Charlier}. This is the pointwise form of
Corollary~\ref{spf:cor:distribution}, which is the stronger statement (it controls
the summed $\ell^1$ error); it asserts no relative-error expansion
uniformly for $j$ growing with $n$. Manuscript~03 makes the first order
explicit for the two sequences: from $B_1^{(1)}=-3u$ and
$B_1^{(2)}=-6u-2u^2$ at $(r,s)=(2,2)$,
\begin{align}
 \Pp(X^{(1)}_{2,2,n}=j)&=\frac{e^{-1}}{j!}
 \Bigl\{1+\frac{3(1-j)}n+\bigO_j(n^{-2})\Bigr\},
 \label{fga:eq:localdirect}\\
 \Pp(X^{(2)}_{2,2,n}=j)&=e^{-2}\frac{2^j}{j!}
 \Bigl\{1+\frac{4-j(j+1)/2}n+\bigO_j(n^{-2})\Bigr\}.
 \label{fga:eq:localabsolute}
\end{align}
At $j=0$ these are the first avoidance corrections; at $j=1,2$ they show
how probability mass is redistributed, which an avoidance count alone does
not.

\section{The inverse in a coordinate absorbing the constant}
\label{fga:sec:inverse}

For A189281 and A110128 the logarithmic coefficients \eqref{spf:eq:a-coeff}
are
\begin{center}
\begin{tabular}{lrrr}
\toprule
Sequence & $a_1$ & $a_2$ & $a_3$\\
\midrule
A189281 & $37/12$ & $-5/2$ & $1439/360$\\
A110128 & $49/12$ & $0$ & $4319/360$\\
\bottomrule
\end{tabular}
\end{center}
The first row and $a_1$ for A110128 are in \cref{spf:sec:inverse}; $a_2$ and
$a_3$ for A110128 are added by manuscript~03.

Manuscript~03 inverts \eqref{spf:eq:logA} in coordinates that absorb
$\kappa_\theta$. For large $Y$ put
\begin{equation}
 \Lambda=\log Y-\kappa_\theta,\qquad N=\frac\Lambda{W_0(\Lambda/e)}
 =\exp\bigl(W_0(\Lambda/e)+1\bigr),\qquad\Delta=\log N,
 \label{fga:eq:coords}
\end{equation}
so that $N(\log N-1)=\Lambda$ exactly. Part~I's coordinates
\eqref{spf:eq:Lambert} are the same construction with $\log Y$ in place of
$\Lambda$; so $N(Y)$ is Part~I's $z$ evaluated at $e^{-\kappa_\theta}Y$.

\begin{proposition}[First inverse corrections; manuscript~03]
\label[proposition]{fga:prop:inverse}
The formal inverse $n_*(Y)$ of \eqref{spf:eq:logA} satisfies
\begin{equation}
 n_*(Y)=N-\frac12+\frac{1/8-a_1}{N\Delta}
 +\frac{1/24-a_1/2-a_2}{N^2\Delta}+\bigO\Bigl(\frac1{N^3\Delta}\Bigr),
 \label{fga:eq:inverse}
\end{equation}
and to every fixed order
$n_*(Y)=N-\frac12+\sum_{j\le J}\omega_j(\Delta)N^{-j}+\bigO(N^{-J-1}\Delta^{-1})$,
with each $\omega_j$ rational in $\Delta$ and computable from
$a_1,\ldots,a_j$; for instance
\begin{equation}
 \omega_3(\Delta)=-\frac1\Delta\Bigl\{\frac{(1/8-a_1)^2}\Delta
 +\frac{(1/8-a_1)^2}{2\Delta^2}+\frac{a_1}4+a_2+a_3-\frac1{64}\Bigr\}.
 \label{fga:eq:omega3}
\end{equation}
For the two sequences,
\begin{equation}
 n_*(Y)=N-\frac12-\frac{71}{24N\Delta}+\frac1{N^2\Delta}+\cdots,
 \qquad
 n_*(Y)=N-\frac12-\frac{95}{24N\Delta}-\frac2{N^2\Delta}+\cdots,
 \label{fga:eq:inverse-ab}
\end{equation}
with $\kappa_1$ and $\kappa_2$ respectively in \eqref{fga:eq:coords}.
\end{proposition}
\begin{proof}[Proof (manuscript~03)]
Put $n=N+\delta$. The dominant part $x\log x-x$ has derivative $\Delta$
and second derivative $1/N$ at $N$; after subtracting $\Lambda$ the
equation becomes
$0=\Delta\delta+\frac{\delta^2}{2N}-\frac{\delta^3}{6N^2}+\cdots+\frac\Delta2
+\frac12\log(1+\delta/N)+\sum_ja_j(N+\delta)^{-j}$. The constant term gives
$\delta=-\frac12+\bigO(1/(N\Delta))$; with $\delta=-\frac12+\sum_j\omega_jN^{-j}$
the next two orders read $\Delta\omega_1+a_1-\frac18=0$ and
$\Delta\omega_2+\frac12a_1+a_2-\frac1{24}=0$. At each later order the new
coefficient enters only as $\Delta\omega_j$. For a fixed truncation the
uncancelled residual is $\bigO(N^{-J-1})$ and the derivative of the smooth
carrier is $\Delta+\bigO(1/N)$, so the mean value theorem divides it by a
quantity asymptotic to $\Delta$.
\end{proof}

The first coefficient agrees with Part~I: $1/8-a_1=1/24-c_1$ is the
coefficient $-(c_1-1/24)$ of \cref{spf:thm:inverse-first}, which in
Part~I's coordinates also carries the terms in $\kappa_\theta$ that are
absorbed here. Manuscript~03 interprets $n_*$ as the inverse of the same
smooth carrier as Part~I, with $X_M(A(n))=n+\bigO(n^{-M-1}/\log n)$, and
brackets the integer threshold $\nu(Y)=\min\{n:A(n)\ge Y\}$ between
$\lceil X_M(Y)\mp\epsilon_Y\rceil$, $\epsilon_Y=\bigO(N^{-M-1}/\log N)$, with
an existence constant only; this bracket is \texttt{p0:thm:staircase} of
\cite{tai}. It notes that expanding $W_0$ in logarithms would give the
familiar nested logarithmic transseries in $Y$, and that $\Lambda$ depends on
$\theta$, so one $\Lambda$ must not be reused for both sequences.

\section{Manuscript 03's computations}
\label{fga:sec:numerics}

Manuscript~03's \texttt{verify.py} (shipped as
\texttt{code/02-fixed-gap-verify.py}) checks, with assertions: 90 avoidance
counts and 90 complete histograms against brute-force permutations for
$n=0,\ldots,8$ and the offset pairs $(1,1),(1,2),(2,2),(2,3),(3,3)$, both
orientations; its exact tiling enumerator against A189281 for
$n=0,\ldots,26$ and $n=40$; 132 stable-profile identities and a
nonrealizable zero case; the moment agreement and disagreement of
\cref{fga:rem:n12}; 50 one-path identities \eqref{fga:eq:onepathmoment};
all posted coefficients of both sequences through order ten, and orders 11
and~12; the general first correction and the mean through order three; and
that the residual of \eqref{fga:eq:inverse-ab} vanishes through $N^{-3}$.
Its exact enumerator and its coefficient engine are independent of each
other and of Part~I's programs. At order twelve the engine enumerates the
77 long profiles of defect twelve, 272 of defect at most twelve, with the
number of dimers summed exactly.

\begin{table}[htbp]
\centering\small
\begin{tabular}{rrrrr}
\toprule
$n$ & $\varepsilon_2(n)$ & $\varepsilon_5(n)$ & $\varepsilon_{10}(n)$ & $\varepsilon_{12}(n)$\\
\midrule
20  & $1.2642\cdot10^{-4}$ & $4.8572\cdot10^{-7}$ & $-2.3683\cdot10^{-10}$ & $5.4066\cdot10^{-11}$\\
40  & $1.5661\cdot10^{-5}$ & $6.9763\cdot10^{-9}$ & $-1.3658\cdot10^{-13}$ & $4.8917\cdot10^{-15}$\\
80  & $1.9541\cdot10^{-6}$ & $1.0406\cdot10^{-10}$ & $-6.8418\cdot10^{-17}$ & $5.2425\cdot10^{-19}$\\
120 & $5.7883\cdot10^{-7}$ & $8.9918\cdot10^{-12}$ & $-7.9394\cdot10^{-19}$ & $2.5777\cdot10^{-21}$\\
\bottomrule
\end{tabular}
\caption{Manuscript~03: the unscaled errors
$\varepsilon_M(n)=e\,\Aone(n)/n!-\sum_{j\le M}c_j^{(1)}(2,2)n^{-j}$, at
110-digit working precision, rounded. The values $\Aone(80)$ and
$\Aone(120)$ are b-file inputs~\cite{fgaBfile}, not regenerated. These are
uncertified decimals; Part~I's \cref{spf:tab:certificates} gives certified
intervals for the scaled errors at $n=80,160,320$.}
\label{fga:tab:errors}
\end{table}

[Added 1 October 2026, batch 74: scaled by $n^{M+1}$, the last two columns
are consistent with Part~I's higher coefficients:
$n^{11}\varepsilon_{10}(n)=-58771$ and $-58990$ at $n=80,120$, against
$c_{11}^{(1)}(2,2)=-59103$; $n^{13}\varepsilon_{12}(n)=2.882\cdot10^6$ and
$2.758\cdot10^6$, approaching $c_{13}^{(1)}(2,2)=2525459$ from above as the
next term $c_{14}^{(1)}/n$ with $c_{14}^{(1)}=26632378$ predicts
($2525459+26632378/80\approx2.858\cdot10^6$).] The four-term inverse
\eqref{fga:eq:inverse-ab} at $Y=\Aone(n)$ misses $n$ by about
$2.0566\cdot10^{-4}$, $1.9555\cdot10^{-5}$, $1.9282\cdot10^{-6}$ and
$5.0462\cdot10^{-7}$ at $n=20,40,80,120$; the signs are not predicted by
the theorem. Manuscript~03 records its run as ``Completed in 22.44
seconds'' and its README speaks of ``tens of seconds''. On this
repository's machine the order-twelve run took more than 170~seconds at the
intake and 104~seconds at the write, and an order-nine run 72~seconds.
The run time depends on the machine and its load.

\section{Further questions raised by manuscript 03}
\label{fga:sec:questions}

Manuscript~03 poses twelve research questions, as proposals only. Six
overlap \cref{spf:sec:questions}: a certificate for the particular
recurrence (Part~I's question~2); integrality beyond the one-path case,
which asks whether $B_J^{(1)}(u;r,s)\in\Z[u]$ for all $r,s$ and so whether
all $c_J^{(1)}(2,2)$ are integers (question~1, and \cref{fga:cor:integrality}
for what is proved); large-order growth (question~4); growing gaps and
component counts (question~5); cycles and other sparse constraint graphs
(question~6); and a small formal package, starting from the profile
identity, the stable polynomial, a verified coefficient checker and the
universality bound (question~7). The other six are new here.
\begin{enumerate}
\item \emph{A sharper oriented degree bound.} Does
$\deg_uB_J^{(1)}(u;r,s)\le J$ hold at every order, sharpening $2J$? The
explicit polynomials through order three show the cancellation; such a
bound would constrain the factorial cumulants of $X$.
\item \emph{Sharp short-path thresholds.} Condition
\eqref{fga:eq:weaklength} is sufficient, not shown optimal. What is the
first shape-dependent term when a shortest path has fixed length, or grows
like $c\log n/\log\log n$?
\item \emph{Optimal quantitative universality.} Find the correct constant
in the $n\log n$ exponent of \cref{fga:thm:universal} for given
proportions, and the path balances that maximize the difference.
\item \emph{Several gaps simultaneously.} The joint law of violations for
several position/value gap pairs, with interacting edge colours and
components containing cycles; a first target is two pairs sharing one gap.
\item \emph{Growing local counts and moderate deviations.} Extend
\eqref{fga:eq:localall} to a quantified range $j=j(n)\to\infty$.
\item \emph{Faster generation and certified errors.} Reduce the cost of
many orders through profile symmetries, falling-factorial bases or direct
Touchard operations, and extract explicit constants from
\cref{fga:thm:allorders} to obtain a certified finite-input threshold
algorithm (compare Part~I's question~8).
\end{enumerate}

\paragraph{OEIS and formalization, as manuscript~03 proposes them.} It
suggests an OEIS addition stating that the expansions through $n^{-10}$
have a recurrence-independent proof, appending orders 11 and~12, and
keeping the recurrences' conjectural status; it submitted nothing, and
Part~I already has orders 11 to~16. It proposes a formalization in three
layers (finite combinatorics; algebra of the stable polynomial and
coefficient generation; analysis), with a coefficient checker that
verifies equality with the finite formula rather than with a guessed
recurrence. It supplies no Lean source, and nothing in Part~II is
formalized.

\paragraph{The shipped files of manuscript~03.}
The ten files of manuscript~03 carry the prefix \texttt{02-fixed-gap-}:
its two scripts in \texttt{code/}, its six recorded outputs and
\texttt{requirements.txt} in \texttt{data/}, and \texttt{sources.md} at the
report root, all byte-identical to the delivery. \texttt{verify.py}
imports its sibling by the delivered name (\texttt{from fixed\_gap import
\ldots}) and writes unprefixed files into \texttt{data/} beside its own
directory; run from \texttt{code/} it would fail to import, and if made to
run it would write beside Part~I's files under unprefixed names. The report
README gives a copy-based recipe.

\addcontentsline{toc}{part}{Appendices}

\appendix
\section{Proof dependencies and archive guide}

The main theorem uses the following chain, all proved in the body of the
article:
\begin{equation*}
 \begin{aligned}
 \text{selected-edge tilings}
 &\Longrightarrow \text{stable polynomial moments}\\
 &\Longrightarrow \text{uniform profile expansion}\\
 &\Longrightarrow \text{finite coefficient formula}.
 \end{aligned}
\end{equation*}
The separate factorial-moment bound justifies discarding the tail and also
proves the total-variation universality theorem. Stirling's expansion and the
definition of Lambert $W$ are standard external analytic inputs for the inverse
section. The recurrence conjectures, the rational-collapse conjecture, and the
2026 holonomicity report are not dependencies of the all-orders theorem.

The archive contains the \LaTeX{} source and compiled PDF, the core exact
coefficient/enumeration library, a validation program, the rational-certificate
program, two optional exploratory programs, and their outputs. The README gives
commands to reproduce the computations and compile the article. The symbolic
collapse script asserts only its explicitly bounded range of identities; its
success is not presented as an infinite proof.
[Added 1 October 2026, batch 73O2: in the shipped report the compiled PDF is
a new build of this text, and the archive's README is replaced by the report
README; \cref{spf:app:provenance} lists what is shipped.]

\section{Provenance and editorial notes}
\label{spf:app:provenance}

[Added 1 October 2026, batch 73O2.]

\paragraph{Source.}
The report was built from one manuscript, manuscript~60 of batch~73 of
ProveIt's incoming-reports intake: the archive \texttt{oeis\_path\_forests.zip}
(main file \texttt{article.tex}, a 21-page PDF titled \emph{Unconditional
all-orders asymptotics for two OEIS permutation sequences}). It arrived in
commit \texttt{aa43cc555} and was placed in commit \texttt{6e193dd4f}, which
deleted the archive from \texttt{docs/incoming}; it survives in
\texttt{aa43cc555}. The manuscript pins no repository revision. It cites the
repository only by the path of the Fabius completion audit~\cite{proveit},
which it uses as an editorial example and calls ``not a dependency''. This
report does not continue that audit, and nothing was filed in the Fabius tree.
[Added 1 October 2026, batch 74: since batch~74 the report is built from
two manuscripts; Part~II and the paragraph ``Part~II (batch~74)'' below
describe the second.]

\paragraph{What the write changed.}
\begin{itemize}
\item Every label gained the prefix \texttt{spf:}; no label was renamed
 otherwise and none was removed. The manuscript had 78 labels.
\item Added: the editorial note after the status paragraph;
 Remark~\ref{spf:rem:offsets-one-one} on the offsets $(1,1)$ (A000255 and
 A002464); Remark~\ref{spf:rem:inverse-apparatus} on the repository's inversion
 apparatus; bracketed notes on Padhi's preprint, on integer thresholds and on
 the shipped files; three bibliography entries~\cite{oeis000255,oeis002464,tai};
 and this section. A separate reciprocal note, Remark~\ref{spf:rem:clique-analogue},
 points to the local law of the companion report on balanced Smirnov words.
\item Nothing was removed. Every non-claim of the manuscript is kept: the status
 paragraph, the Poincar\'e-only reading of Table~\ref{spf:tab:extra}, the
 limitation after \cref{spf:thm:universality}, the fixed-offset limitation
 after the proof of \cref{spf:thm:main}, the conditional nature of
 Proposition~\ref{spf:prop:collapse-implication}, and the statement that the
 recurrences of A189281 and A110128 are not proved.
\end{itemize}

\paragraph{Status of the conjecture.}
Conjecture~\ref{spf:conj:collapse} (rational collapse) is a conjecture. It has been
checked as an identity of rational functions for $h=4,\ldots,16$ only, and the
integrality of every $c_J^{(1)}(2,2)$ is conditional on it
(Proposition~\ref{spf:prop:collapse-implication}).

\paragraph{Overloaded letters.}
The manuscript reuses a few letters; none is renamed here. $C$ is a tiling
count $C_{\mathcal P}(\mathbf b)$ in \cref{spf:sec:tilings,spf:sec:stable}, the
number of non-dimer components in \cref{spf:sec:coefficients}, a constant
$C_{r,s}$ in Lemma~\ref{spf:lem:Qbound}, and the series $C(x)$ in
\cref{spf:sec:inverse}. $K$ is a moment cutoff in
Lemma~\ref{spf:lem:moment-bound} and Theorem~\ref{spf:thm:universality} but the edge degree of the
non-dimer part of a profile in \cref{spf:sec:coefficients}. $d$ is the
defect of a profile, while $d_j(w)$ are the inverse coefficients of
\cref{spf:sec:inverse}. $\mathcal D$ is a formal series in
\cref{spf:sec:stable} and $D$ the index correction in \cref{spf:sec:inverse}.
In each case the section determines the meaning.

\paragraph{Checks by the intake.}
The intake reran every shipped program on a copy (all passed; the
regenerated outputs equal the shipped ones apart from line endings, a final
newline and timings), recounted both sequences for $n\le9$ by brute force,
compared the two order-ten expansions with the OEIS entries digit for digit,
rederived $d_1$ in \cref{spf:thm:inverse-first} by hand, and made the
$(1,1)$ checks of Remark~\ref{spf:rem:offsets-one-one}. These checks do not
replace an independent review of the proofs.

\paragraph{Part~II (batch 74).}
[Added 1 October 2026, batch 74.] Part~II is built from manuscript~03 of
batch~74, the archive \nolinkurl{ProveIt_Fixed_Gap_Asymptotics.zip}
(14~files, a 25-page PDF), which arrived in commit \texttt{c664fc2f0} and
was placed in commit \texttt{b669cff87}. Its pinned commit is
\texttt{291fbe6bb}~\cite{fgaProveIt}. Its article, README, PDF and checksum
ledger \texttt{SHA256SUMS.txt} (13/13 verified at placement) are not
shipped; its other ten files are, byte-identical, with the prefix
\texttt{02-fixed-gap-}: \texttt{code/fixed\_gap.py} and
\texttt{code/verify.py} in \texttt{code/}; the six files of its
\texttt{data/} and \texttt{requirements.txt} in \texttt{data/}; and
\texttt{sources.md} at the report root. Those files use the delivery names.

\paragraph{Where the batch-74 merge had to choose.} Manuscript~60 stays the
base, being the more general in $u$ and the stronger in arithmetic
(denominator theorem, order sixteen), in the whole-law statement and in
its certificates. Manuscript~03's re-derivations of Part~I's theorems are
recorded as second routes in \cref{fga:tab:routes} and not reprinted; its
own results form Part~II, printed after Part~I's conclusion so that no
existing number changes. Its symbols are renamed into Part~I's notation
(\cref{fga:tab:dictionary}); the devices its new proofs need (the moment
bound with $D_{n,k}$, the unit-disk Taylor bound, the defect-sensitive
bound) are printed with them. Its coefficients through order twelve are
printed only as a confirmation of \cref{spf:tab:extra}.

\paragraph{What the batch-74 write changed.} It added the part headings,
the editorial note after the batch-73O2 one, bracketed notes marked
``[Added 1 October 2026, batch 74]'' after the proof of \cref{spf:thm:main},
in questions~1 and~5 of \cref{spf:sec:questions} and in this appendix,
Part~II (\cref{fga:sec:front} to \cref{fga:sec:questions}), two bibliography
entries, and a table-of-contents line for the appendices. All its labels
carry the prefix \texttt{fga:}. No existing label was renamed or removed,
no existing number changed, and no delivered sentence of manuscript~60 was
changed or removed. Manuscript~03's text is restated in Part~I's notation,
with some proofs condensed; none of its statements was strengthened.

\begin{thebibliography}{99}
\bibitem{oeis189281}
OEIS Foundation Inc., \emph{A189281: Number of permutations $p$ of
$1,2,\ldots,n$ satisfying $p(i+2)-p(i)\ne2$}.
\url{https://oeis.org/A189281}. Entry inspected 1 October 2026.

\bibitem{oeis110128}
OEIS Foundation Inc., \emph{A110128: Number of permutations $p$ of
$1,2,\ldots,n$ satisfying $|p(i+2)-p(i)|\ne2$}.
\url{https://oeis.org/A110128}. Entry inspected 1 October 2026.

\bibitem{spahn}
G. Spahn and D. Zeilberger,
\emph{Counting permutations where the difference between entries located
$r$ places apart can never be $s$ (for any given positive integers $r$ and $s$)},
arXiv:2211.02550 (2022), revised manuscript with postscript.
\url{https://sites.math.rutgers.edu/~zeilberg/mamarim/mamarimPDF/permsV2.pdf}.

\bibitem{tauraso}
R. Tauraso, \emph{The dinner table problem: the rectangular case},
Integers \textbf{6} (2006), A11; arXiv:math/0507293.
\url{https://arxiv.org/abs/math/0507293}.

\bibitem{kauers}
M. Kauers and C. Koutschan, \emph{Guessing with little data},
Proceedings of ISSAC 2022, 83--90; arXiv:2202.07966.
\url{https://arxiv.org/abs/2202.07966}.

\bibitem{padhi}
J. S. Padhi, \emph{Solutions to five challenge problems in enumerative and
algorithmic combinatorics, with an account of the human--machine methodology
employed}, arXiv:2608.11290 (2026), especially Sections 5 and 6.
\url{https://arxiv.org/pdf/2608.11290}. Consulted as a report of current
related work; not used as a theorem dependency.

\bibitem{dlmfGamma}
NIST Digital Library of Mathematical Functions,
\emph{Section 5.11: Asymptotic expansions of the gamma function}.
\url{https://dlmf.nist.gov/5.11}. Inspected 1 October 2026.

\bibitem{dlmfW}
NIST Digital Library of Mathematical Functions,
\emph{Section 4.13: Lambert $W$-function}.
\url{https://dlmf.nist.gov/4.13}. Inspected 1 October 2026.

\bibitem{proveit}
V. Reshetnikov, \emph{ProveIt}, repository and asymptotic completion audit:\\
\path{Analysis/FabiusFunction/docs/ASYMPTOTIC_COMPLETION_AUDIT.md}.
\url{https://github.com/VladimirReshetnikov/ProveIt}.
Repository material inspected 1 October 2026. Used as project context and an
example of separating asymptotic assertions from their verification status.

\bibitem{oeis000255}
[Added in the batch-73O2 write.] OEIS Foundation Inc., \emph{A000255:
$a(n)=n\,a(n-1)+(n-1)\,a(n-2)$, $a(0)=a(1)=1$}.
\url{https://oeis.org/A000255}. Entry inspected 1 October 2026.

\bibitem{oeis002464}
[Added in the batch-73O2 write.] OEIS Foundation Inc., \emph{A002464:
Hertzsprung's problem}, including the Abramson--Moser asymptotic expansion
computed by V.~Kotesovec. \url{https://oeis.org/A002464}. Entry inspected
1 October 2026.

\bibitem{tai}\begingroup\raggedright
[Added in the batch-73O2 write.] V. Reshetnikov and contributors,
\emph{Transseries and Inversion}, canonical volume of the ProveIt
series-and-transseries collection. Its source is
\path{Analysis/Transseries/docs/series-and-transseries/}\allowbreak
\path{Transseries_And_Inversion/transseries_and_inversion.tex}. Cited here:
Part ``Inversion: the apparatus for a rapidly growing function''
(labels \texttt{p0:thm:perturbed-inversion}, \texttt{p0:thm:staircase}),
chapter on the increasing inverse of Gamma (\texttt{p6:thm:gamma}) and chapter
``The subfactorial'' (\texttt{p8:sec:top}). Repository state of placement
commit \texttt{6e193dd4f}.\par\endgroup

\bibitem{fgaBfile}
[Added in the batch-74 write, from manuscript 03.] R. Matsuo, V. Kotesovec
and C. Koutschan, \emph{Table of A189281 for $n=0,\ldots,300$}, with the
authorship of ranges given in the OEIS entry.
\url{https://oeis.org/A189281/b189281.txt}. Accessed by manuscript~03 on
1 October 2026.

\bibitem{fgaProveIt}\begingroup\raggedright
[Added in the batch-74 write, from manuscript 03.] V. Reshetnikov,
\emph{ProveIt}, repository and README, inspected 1 October 2026 at commit
\nolinkurl{291fbe6bb4477e0cb2131f1e041b12a9abe363fd}.
\url{https://github.com/VladimirReshetnikov/ProveIt}.\par\endgroup
\end{thebibliography}
\end{document}
